\chapter{Vector Fields}

\section{Vector Fields Introduction}

Let's begin the vector field discussion by plotting a simple vector field $\mathbf{F} = \left< -y,x \right>$ in SageMath.  It field verctors are rotating around the origin.

\begin{programcode}{SageMath Code and Output:}
\begin{verbatim}
var('x,y,z')
plot_vector_field([-y,x], (x,-5,5), (y,-5,5))
\end{verbatim}
\end{programcode}
The output is shown in Fig \ref{FieldRotate}
\begin{figure}[H]
\sidecaption[t]
\includegraphics[scale=.65]{chapters/ch10VectorFields/FieldRotate}
\caption{Vector field rotating}
\label{FieldRotate}       % Give a unique label
\end{figure}

Next, let's consider a conservative field $\mathbf{F} = \left< 2x,4y \right>$, which is the gradient of $f(x,y)=x^2+2y^2$.  We also check the conservative field condition for a general field $\mathbf{F} = \left< M,N \right>$ as $\frac{\partial M}{\partial y} = 0 = \frac{\partial N}{\partial x}$.

\begin{programcode}{SageMath Code and Output:}
\begin{verbatim}
reset(); var('x,y')
f(x,y) = x^2 + 2*y^2 
F(x,y) = f(x,y).gradient()
show('f(x,y) = ',f(x,y),'  F(x,y) = ',F(x,y))
M(x,y) = F(x,y)[0];   N(x,y) = F(x,y)[1];
show('M(x,y) = ', M(x,y), '  N(x,y) = ', N(x,y))
show('dM/dy = ', M(x,y).diff(y),
     '    dN/dx = ', N(x,y).diff(x))
plot_vector_field(F(x,y), (x,-5,5), (y,-5,5), 
                  figsize = (4,3))
\end{verbatim}
\end{programcode}
The output is shown in Fig \ref{ConservField}
\begin{figure}[H]
\sidecaption[t]
\includegraphics[scale=.65]{chapters/ch10VectorFields/ConservField}
\caption{Conservative vector field}
\label{ConservField}       % Give a unique label
\end{figure}

Let's consider another conservative field $\mathbf{F} = \left< 3+2xy,x^2 - 3y^2 \right>$. The SageMath cell below checks the conservative field condition and computes antiderivatives of $M$ and $N$ components to construct $f$:  \
$\int F_x dx = \int 3+2xy dx = 3x+x^2y $,  $\int F_y dy = \int x^2 - 3y^2 dy = x^2y-y^3 $,  which results in potential function $f(x,y) =  x^2y+3x-y^3$. 
Then, we plot level curves of $f$ and vector field $\mathbf{F}$ together, confirming that they are perpendicular to each other.

\begin{programcode}{SageMath Code and Output:}
\begin{verbatim}
reset(); var('x,y')
F(x,y) = [3+2*x*y, x^2 - 3*y^2]
show('int Fx in x = ', integrate(F(x,y)[0],x),
   '  int Fy in y = ', integrate(F(x,y)[1],y))

M(x,y) = F(x,y)[0];   N(x,y) = F(x,y)[1];
show('M(x,y) = ', M(x,y), '  N(x,y) = ', N(x,y))
show('dM/dy = ', M(x,y).diff(y),
     '    dN/dx = ', N(x,y).diff(x))

show('Based on antiderivatives above we obtain:')
f(x,y) = x^2*y-y^3+3*x
show('f(x,y) = ', f(x,y))
p1 = contour_plot(f(x,y), (x,-5,5), (y,-5,5), figsize = (4,3))
p2 = plot_vector_field(F(x,y), (x,-5,5), (y,-5,5))
(p1+p2).show()
\end{verbatim}
\end{programcode}
The output is shown in Fig \ref{ConservField1}
\begin{figure}[H]
\sidecaption[t]
\includegraphics[scale=.65]{chapters/ch10VectorFields/ConservField1}
\caption{Conservative vector field}
\label{ConservField1}       % Give a unique label
\end{figure}

In 3D, the $curl(\mathbf{F})=0$ is the condition for a vector field to be conservative.   For example, vector field $\mathbf{F} = \left< 2xy,x^2+2yz,y^2 \right>$ is conservative:
$$curl(\mathbf{F}) = \begin{pmatrix}
\mathbf{i}& \mathbf{j} & \mathbf{k}\\
\frac{\partial}{\partial x} & \frac{\partial}{\partial y} & \frac{\partial}{\partial z} \\
2xy & x^2+2yz & y^2
\end{pmatrix} =(2y-2y) \mathbf{i} - (0-0) \mathbf{j}+ (2x-2x) \mathbf{k} =   \mathbf{0} $$
Doing partial integration:  $\int F_x dx = \int 2xy dx = x^2y $,  $\int F_y dy = \int x^2+2yz dy = x^2y+y^2z $, $\int F_z dz = \int y^2 dz = y^2z $ results in potential function $f(x,y,z) =  x^2y+y^2z$.   The SageMath cells below check that curl is 0, perform partial integration, and check that $grad(f)=\nabla f = \mathbf{F}$.

\begin{programcode}{SageMath Code and Output:}
\begin{verbatim}
reset(); var('x,y,z')
from sage.all import *
F(x,y,z) = vector([2*x*y,x^2+2*y*z,y^2])
show('F = ',F(x,y,z))
curlF(x,y,z) = F.curl()(x,y,z)
show('curl of F = ', curlF(x,y,z))
show('int Fx in x = ', integrate(F(x,y,z)[0],x))
show('int Fy in y = ', integrate(F(x,y,z)[1],y))
show('int Fz in z = ', integrate(F(x,y,z)[2],z))

f(x,y,z) = x^2*y+y^2*z
show('f(x,y,z) = ', f(x,y,z))
show(' grad of f = ',f.gradient()(x,y,z))

plot_vector_field3d([2*x*y,x^2+2*y*z,y^2], 
   (x,-5,5), (y,-5,5),(z,-5,5),thickness=5)
\end{verbatim}
\end{programcode}
The output is shown in Fig \ref{Conservative3D}
\begin{figure}[H]
\sidecaption[t]
\includegraphics[scale=.65]{chapters/ch10VectorFields/Conservative3D}
\caption{Conservative 3D field}
\label{Conservative3D}       % Give a unique label
\end{figure}

As an example of a 3D field which is not conservative let's consider $\mathbf{F} = \left< \frac{y}{x},-\frac{x^2}{y^2},2xz \right>$.  The curl computed below is not a zero vector.
$$curl(\mathbf{F}) = \begin{pmatrix}
\mathbf{i}& \mathbf{j} & \mathbf{k}\\
\frac{\partial}{\partial x} & \frac{\partial}{\partial y} & \frac{\partial}{\partial z} \\
\frac{y}{x} & -\frac{x^2}{y^2} & 2xz
\end{pmatrix} =(0-0) \mathbf{i} - (2z-0) \mathbf{j}+ (-\frac{2x}{y}-\frac{1}{x}) \mathbf{k} $$

\begin{programcode}{SageMath Code and Output:}
\begin{verbatim}
reset(); var('x,y,z')
F(x,y,z) = vector([y/x,-x^2/y^2,2*x*z])
show('F = ',F(x,y,z))
curlF(x,y,z) = F.curl()(x,y,z)
show('curl of F = ', curlF(x,y,z))
\end{verbatim}
\end{programcode}
The output is shown in Fig \ref{NonConservative3D}
\begin{figure}[H]
\sidecaption[t]
\includegraphics[scale=.65]{chapters/ch10VectorFields/NonConservative3D}
\caption{A non-conservative 3D vector field}
\label{NonConservative3D}       % Give a unique label
\end{figure}

It can be shown that the inverse square field is conservative for the potential function below:
$$ f(x,y,z) = \frac{-K}{\sqrt{x^2+y^2+z^2}} \ \ \Rightarrow \ \ f_x = \frac{Kx}{(x^2+y^2+z^2)^{3/2}}$$
By symmetry, other derivatives are very similar, so 
$$ \nabla f = \left<   \frac{Kx}{(x^2+y^2+z^2)^{3/2}},  \frac{Ky}{(x^2+y^2+z^2)^{3/2}} ,  \frac{Kz}{(x^2+y^2+z^2)^{3/2}}  \right> = K \frac{\mathbf{r}}{||\mathbf{r}||^2} \frac{\mathbf{r}}{||\mathbf{r}||}$$
If we let $\mathbf{F} = \nabla f$, then $curl(\mathbf{F}) = \nabla \times (\nabla f) =\mathbf{0} $.
Note that we must exclude the origin from plotting since the field becomes infinite at that point.

\begin{programcode}{SageMath Code and Output:}
\begin{verbatim}
reset();  var('x,y,z,K')
f(x,y,z) = -K/sqrt(x^2+y^2+z^2)
fx(x,y,z) = derivative(f(x,y,z),x)
fy(x,y,z) = derivative(f(x,y,z),y)
fz(x,y,z) = derivative(f(x,y,z),z)
show('f(x,y,z) = ',f(x,y,z))
show('fx(x,y,z) = ',fx(x,y,z))
show('fy(x,y,z) = ',fy(x,y,z))
show('fz(x,y,z) = ',fz(x,y,z))
F(x,y,z) = vector([fx(x,y,z),fy(x,y,z),fz(x,y,z)])
curlF(x,y,z) = F.curl()(x,y,z)
show('curl of F = ', curlF(x,y,z))

d = (x^2+y^2+z^2)^(3/2)
plot_vector_field3d([x/d,y/d,z/d], 
   (x,0.5,2), (y,0.5,2),(z,0.5,2),thickness=5)
\end{verbatim}
\end{programcode}
The output is shown in Fig \ref{InverseSquare}
\begin{figure}[H]
\sidecaption[t]
\includegraphics[scale=.65]{chapters/ch10VectorFields/InverseSquare}
\caption{Inverse square field}
\label{InverseSquare}       % Give a unique label
\end{figure}

As an example of divergence of a 3D vector field, let's consider vector field $\mathbf{F} = \left< M,N,P \right> = \left< x^2y^3, y^2z, xz^2 \right>$.
The divergence is $div(\mathbf{F})=\frac{\partial M}{\partial x} + \frac{\partial N}{\partial y} + \frac{\partial P}{\partial z} = 2xy^3+2yz+2xz$.

\begin{programcode}{SageMath Code and Output:}
\begin{verbatim}
reset(); var('x,y,z')
F(x,y,z) = vector([x^2*y^3,y^2*z,x*z^2])
divF(x,y,z) = F.div()(x,y,z)
show('F = ',F(x,y,z),'   div of F = ', divF(x,y,z))
\end{verbatim}
\end{programcode}
The output is shown in Fig \ref{Divergence3D}
\begin{figure}[H]
\sidecaption[t]
\includegraphics[scale=.65]{chapters/ch10VectorFields/Divergence3D}
\caption{Divergence example}
\label{Divergence3D}       % Give a unique label
\end{figure}








\section{Line Integrals}

Let's begin line integrals consideration with the simplest example of integrating function $f(x,y,z) =2x^2-3y^2+4z$
on a line from the origin $O(0,0,0)$ to point $P(2,1,3)$.  The equation of the line segment $OP$ is $x = 0+2t=2t; \ \ y=0+t=t;
 \ \ z = 0+3t=3t \ \ \ 0 \le t \le 1$.  Therefore, the length element is $ds=\sqrt{x'(t)^2+y'(t)^2+z'(t)^2}dt=\sqrt{2^2+1^2+3^2}dt=\sqrt{14}dt$.
$$\int_L fds = \int_0^1 (2(2t)^2-3(t)^2+4(3t) )\sqrt{14}dt = \sqrt{14} \int_0^1 5t^2+12t dt =  $$
$$ \sqrt{14} \left( 5\frac{t^3}{3}+12\frac{t^2}{2} \right)_0^1 = \sqrt{14} \left( \frac{5}{3} + 6 \right) = \sqrt{14} \frac{23}{3} $$

\begin{programcode}{SageMath Code and Output:}
\begin{verbatim}
reset(); var('x,y,z,t')
f(x,y,z) = 2*x^2-3*y^2+4*z
x1(t)=2*t;  y1(t)=t;  z1(t)=3*t; t0=0; tf=1

x1t(t) = derivative(x1(t),t)
y1t(t) = derivative(y1(t),t)
z1t(t) = derivative(z1(t),t)
dsdt = sqrt(x1t(t)^2+y1t(t)^2+z1t(t)^2)

show('f(x,y,z) = ',f(x,y,z))
show('x1(t) = ',x1(t),' y1(t) = ',y1(t),
     ' z1(t) = ',z1(t))
show('x1t(t) = ',x1t(t),' y1t(t) = ',y1t(t),
     ' z1t(t) = ',z1t(t))
show(' dsdt = ',dsdt)

Integrant = f(x1(t),y1(t),z1(t))*dsdt
show('Integrant = ', Integrant)
Int(t) = integrate( Integrant, t)
show('Int(t) = ', Int(t))
Int1 = Int(tf) - Int(t0)
show('Int1 = ',Int1)

parametric_plot([x1(t),y1(t)],(t,0,2*pi),
                       figsize = (4,3))
\end{verbatim}
\end{programcode}
The output is shown in Fig \ref{LineInt1}
\begin{figure}[H]
\sidecaption[t]
\includegraphics[scale=.65]{chapters/ch10VectorFields/LineInt1}
\caption{Line integral example}
\label{LineInt1}       % Give a unique label
\end{figure}



As an application of this type of line integral, let's find the mass of one coil of a helical spring $\mathbf{r}(t) = \left< 3\cos(t), 3\sin(t), \sqrt{7}t \right> \ \ \  0 \le t \le 2\pi$ with density dependent on the height as $\rho=z^2+3$.  The derivative of the path function is  $\mathbf{r}'(t) =   \left< -3\sin(t), 3\cos(t), \sqrt{7} \right> \ \Rightarrow \ ||\mathbf{r'}(t)|| = \sqrt{9\sin^2(t)+9\cos^2(t)+7}=\sqrt{16}=4$, so the length element is $ds =||\mathbf{r'}(t)||dt =4dt $.  Therefore, the mass is:
$$m = \int_L fds = \int_0^{2\pi} ((\sqrt{7}t)^2+3)4dt = 4 \left( 7\frac{t^3}{3}+3t  \right)_0^{2\pi} = 24\pi+\frac{224\pi^3}{3} $$

\begin{programcode}{SageMath Code and Output:}
\begin{verbatim}
reset(); var('x,y,z,t')
f(x,y,z) = z^2+3
x1(t)=3*cos(t);  y1(t)=3*sin(t); 
z1(t)=sqrt(7)*t; t0=0; tf=2*pi

x1t(t) = derivative(x1(t),t)
y1t(t) = derivative(y1(t),t)
z1t(t) = derivative(z1(t),t)
dsdt = sqrt(x1t(t)^2+y1t(t)^2+z1t(t)^2).full_simplify()

show('f(x,y,z) = ',f(x,y,z))
show('x1(t) = ',x1(t),' y1(t) = ',y1(t),' z1(t) = ',z1(t))
show('x1t(t) = ',x1t(t),' y1t(t) = ',y1t(t),' z1t(t) = ',z1t(t))
show(' dsdt = ',dsdt)

Integrant = f(x1(t),y1(t),z1(t))*dsdt
show('Integrant = ', Integrant)
Int(t) = integrate( f(x1(t),y1(t),z1(t))*dsdt, t)
show('Int(t) = ', Int(t))
Int1 = Int(tf) - Int(t0)
show('Int1 = ',Int1)

parametric_plot3d([x1(t),y1(t),z1(t)],
                  (t,0,2*pi),thickness=5)
\end{verbatim}
\end{programcode}
The output is shown in Fig \ref{LineInt2}
\begin{figure}[H]
\sidecaption[t]
\includegraphics[scale=.65]{chapters/ch10VectorFields/LineInt2}
\caption{Line integral on a spiral}
\label{LineInt2}       % Give a unique label
\end{figure}


There is a different type of line integral for vector fields rather than scalar functions.   One of its applications is to find the work done by a vector field on a given path.  For example, let's find work done by force field $\mathbf{F} = \left< 2x, -3y, 2z \right>$ along the same one coil of a helical spring $\mathbf{r}(t) = \left< 3\cos(t), 3\sin(t), \sqrt{7}t \right> \ \ \  0 \le t \le 2\pi$.  In the same way, $\mathbf{r}'(t) =   \left< -3\sin(t), 3\cos(t), \sqrt{7} \right>$, so the dot product in the integral becomes: 
$$\mathbf{F} \cdot \mathbf{r}'(t) = \left< 2 \cdot 3\cos(t), -3 \cdot 3\sin(t), 2 \sqrt{7}t  \right> \cdot    \left< -3\sin(t), 3\cos(t), \sqrt{7} \right> $$
$$=-18\cos(t)\sin(t)-27\cos(t)\sin(t)+14t= -\frac{45}{2}\sin(2t)+14t$$
Then the work integral is:
$$W = \int_L \mathbf{F} \cdot \mathbf{dr} = \int_L \mathbf{F} \cdot \mathbf{r'}(t)dt =\int_0^{2\pi}\left( -\frac{45}{2}\sin(2t)+14t\right)dt =  $$
$$\left( \frac{45}{4}\cos(2t)+14 \frac{t^2}{2} \right)_0^{2\pi} = 28\pi^2$$

\begin{programcode}{SageMath Code and Output:}
\begin{verbatim}
reset(); var('x,y,z,t')
F(x,y,z) = vector([2*x,-3*y,2*z])
r(t) = vector([3*cos(t),3*sin(t),sqrt(7)*t])
t0=0; tf=2*pi
rp(x,y,z) = derivative(r(t),t)

show('F(x,y,z) = ',F(x,y,z))
show('r(t) = ',r(t))
show('rp(t) = ',rp(t))

Fdr(t) = \ 
 F(r(t)[0],r(t)[1],r(t)[2]).dot_product(rp(t))
show('Fdr(t) = ',Fdr(t))
Int(t) = integrate( Fdr(t), t)
show('Int(t) = ', Int(t))
Int1 = Int(tf) - Int(t0)
show('Int1 = ',Int1)
p1 = plot_vector_field3d([2*x,-3*y,2*z], 
  (x,-5,5), (y,-5,5),(z,0,4*pi),thickness=5)
p2 = parametric_plot3d(r(t), 
        (t,0,2*pi),thickness=5)
(p1+p2).show()
\end{verbatim}
\end{programcode}
The output is shown in Fig \ref{LineIntVectorField}
\begin{figure}[H]
\sidecaption[t]
\includegraphics[scale=.65]{chapters/ch10VectorFields/LineIntVectorField}
\caption{Line integral of a vector field on a spiral}
\label{LineIntVectorField}       % Give a unique label
\end{figure}


There is another way to write this type of line integral of a vector field.   Specifically $\int_L \mathbf{F} \cdot d\mathbf{r} = \int_L \left< M,N,P \right> \cdot  \left< dx,dy,dz \right> =\int_L Mdx+Ndy+Pdz$ or in 2D $\int_L \mathbf{F} \cdot d\mathbf{r} =\int_L Mdx+Ndy$.  For example, let's compute $\int_L 2ydx+3xdy$ over a circle $x=\cos(t); \ \ y=\sin(t); \ \ 0 \le t \le 2\pi$.
$$ \int_L 2ydx+3xdy = \int_0^{2\pi}  2\sin(t)(-\sin(t))dt +3\cos(t)\cos(t)dt =$$
$$ \int_0^{2\pi}  -2+5\cos^2(t)dt = \int_0^{2\pi}  -2+5\frac{1+\cos(2t)}{2}dt = \left( \frac{1}{2}t + \frac{5}{4}\sin(2t) \right)_0^{2\pi} =\pi$$

\begin{programcode}{SageMath Code and Output:}
\begin{verbatim}
reset(); var('x,y,z,t')
M(x,y) = 2*y;  N(x,y) = 3*x;  t0=0; tf=2*pi
x1(t) = cos(t);  y1(t) = sin(t)
x1p(t) = derivative(x1(t),t); 
y1p(t) = derivative(y1(t),t)

show('M(x,y) = ',M(x,y),'  N(x,y) = ',N(x,y))
show('x1(t) = ',x1(t),'  y1(t) = ',y1(t))
show('x1p(t) = ',x1p(t),'  y1p(t) = ',y1p(t))

Fdr(t) = M(x1(t),y1(t))*x1p(t) + \
         N(x1(t),y1(t))*y1p(t)
show('Fdr(t) = ',Fdr(t))
Int(t) = integrate( Fdr(t), t)
show('Int(t) = ', Int(t))
Int1 = Int(tf) - Int(t0)
show('Int1 = ',Int1)
p1 = plot_vector_field([M(x,y),N(x,y)], 
         (x,-2,2), (y,-2,2),figsize = (4,3))
p2 = parametric_plot([x1(t),y1(t)], \
         (t,0,2*pi),thickness=5)
(p1+p2).show()
\end{verbatim}
\end{programcode}
The output is shown in Fig \ref{LineIntVectField2d}
\begin{figure}[H]
\sidecaption[t]
\includegraphics[scale=.65]{chapters/ch10VectorFields/LineIntVectField2d}
\caption{Line integral of a vector field in 2D with scalar product notation}
\label{LineIntVectField2d}       % Give a unique label
\end{figure}










\section{Conservative Vector Fields and Independence of Path}

Let's begin the independence of path discussion by considering a conservative vector field  $\mathbf{F} = \left< 4xy,2x^2 \right>$, which is the gradient of $f(x,y)=2x^2y$.  Also the conservative field condition check confirms it: $\frac{\partial M}{\partial y} = 4x = \frac{\partial N}{\partial x} $.    In the next three SageMath cells below, we consider three different paths.  First, a simple straight line path $\mathbf{r}(t) = \left< t, t \right> \ \ \ 0\le t \le 1$ from $(0,0)$ to $(1,1)$ is used.  The 2nd path is the parabola  $\mathbf{r}(t) = \left< t, t^2 \right> \ \ \ 0\le t \le 1$ between the same points.   Finally, the square root path is given by  $\mathbf{r}(t) = \left< t, \sqrt{t} \right> \ \ \ 0\le t \le 1$.   In all three cases, the line (path) integral $\int_L \mathbf{F} d\mathbf{r} = 2$.    Analytically, along thestraight line path, $\mathbf{r}'(t) = \left< 1, 1 \right>$, so $\mathbf{F}d\mathbf{r} = \left< 4tt,2t^2 \right>\cdot\left< 1, 1 \right>=6t^2$, therefore $\int_L \mathbf{F}d\mathbf{r} = \int_0^1 6t^2 dt = \left( 2t^3 \right)_0^1=2$.   In SageMath, we obtain:

\begin{programcode}{SageMath Code and Output:}
\begin{verbatim}
reset(); var('x,y,z,t')
F(x,y) = vector([4*x*y,2*x^2])
r(t) = vector([t,t])
t0=0; tf=1
rp(x,y,z) = derivative(r(t),t)

show('F(x,y) = ',F(x,y))
show('r(t) = ',r(t))
show('rp(t) = ',rp(t))

Fdr(t) = F(r(t)[0],r(t)[1]).dot_product(rp(t))
show('Fdr(t) = ',Fdr(t))
Int(t) = integrate( Fdr(t), t)
show('Int(t) = ', Int(t))
Int1 = Int(tf) - Int(t0)
show('Int1 = ',Int1)
p1 = plot_vector_field(F(x,y),(x,0,1),(y,0,1))
p2 = parametric_plot(r(t),(t,t0,tf),figsize=(4,3))
(p1+p2).show()
\end{verbatim}
\end{programcode}
The output is shown in Fig \ref{PathIndepLine}
\begin{figure}[H]
\sidecaption[t]
\includegraphics[scale=.65]{chapters/ch10VectorFields/PathIndepLine}
\caption{Path independence, line integral on straight line}
\label{PathIndepLine}       % Give a unique label
\end{figure}

Next, along the second path, $\mathbf{r}'(t) = \left< 1, 2t \right>$, so $\mathbf{F}d\mathbf{r} = \left< 4tt^2,2t^2 \right>\cdot\left< 1, 2t \right>=8t^3$, therefore $\int_L \mathbf{F}d\mathbf{r} = \int_0^1 8t^3 dt = \left( 2t^4 \right)_0^1=2$.   Analogously, in SageMath:

\begin{programcode}{SageMath Code and Output:}
\begin{verbatim}
reset(); var('x,y,z,t')
F(x,y) = vector([4*x*y,2*x^2])
r(t) = vector([t,t^2])
t0=0; tf=1
rp(x,y,z) = derivative(r(t),t)

show('F(x,y) = ',F(x,y))
show('r(t) = ',r(t))
show('rp(t) = ',rp(t))

Fdr(t) = F(r(t)[0],r(t)[1]).dot_product(rp(t))
show('Fdr(t) = ',Fdr(t))
Int(t) = integrate( Fdr(t), t)
show('Int(t) = ', Int(t))
Int1 = Int(tf) - Int(t0)
show('Int1 = ',Int1)
p1 = plot_vector_field(F(x,y),(x,0,1),(y,0,1))
p2 = parametric_plot(r(t),(t,t0,tf),figsize=(4,3))
(p1+p2).show()
\end{verbatim}
\end{programcode}
The output is shown in Fig \ref{PathIndepParabola}
\begin{figure}[H]
\sidecaption[t]
\includegraphics[scale=.65]{chapters/ch10VectorFields/PathIndepParabola}
\caption{Path independence, line integral on parabola path}
\label{PathIndepParabola}       % Give a unique label
\end{figure}

Finally, along the third path, $\mathbf{r}'(t) = \left< 1, \frac{1}{2\sqrt{t}} \right>$, so $\mathbf{F}d\mathbf{r} = \left< 4t\sqrt{t},2t^2 \right>\cdot\left< 1, \frac{1}{2\sqrt{t}} \right>=5t^{3/2}$, therefore $\int_L \mathbf{F}d\mathbf{r} = \int_0^1 5t^{3/2} dt = \left( 5\frac{t^{5/2}}{5/2} \right)_0^1=2$.   In SageMath:

\begin{programcode}{SageMath Code and Output:}
\begin{verbatim}
reset(); var('x,y,z,t')
F(x,y) = vector([4*x*y,2*x^2])
r(t) = vector([t,sqrt(t)])
t0=0; tf=1
rp(x,y,z) = derivative(r(t),t)

show('F(x,y) = ',F(x,y))
show('r(t) = ',r(t))
show('rp(t) = ',rp(t))

Fdr(t) = F(r(t)[0],r(t)[1]).dot_product(rp(t))
show('Fdr(t) = ',Fdr(t))
Int(t) = integrate( Fdr(t), t)
show('Int(t) = ', Int(t))
Int1 = Int(tf) - Int(t0)
show('Int1 = ',Int1)
p1 = plot_vector_field(F(x,y),(x,0,1),(y,0,1))
p2 = parametric_plot(r(t),(t,t0,tf),figsize=(4,3))
(p1+p2).show()
\end{verbatim}
\end{programcode}
The output is shown in Fig \ref{PathIndepSqrt}
\begin{figure}[H]
\sidecaption[t]
\includegraphics[scale=.65]{chapters/ch10VectorFields/PathIndepSqrt}
\caption{Path independence, line integral on square root path}
\label{PathIndepSqrt}       % Give a unique label
\end{figure}

The path independence for a conservative field is proved in general, and in fact, the path integral is equal to the difference of the potential function at the endpoints:
$$\int_L \mathbf{F}d\mathbf{r} = \int_L \nabla f d\mathbf{r} = f(end)-f(start) = \left( 2x^2 y \right)_{(0,0)}^{(1,1)}=2-0 =2$$
Therefore, for a conservative vector field, a path need not even be specified - only the starting and ending points.   

\begin{programcode}{SageMath Code and Output:}
\begin{verbatim}
reset(); var('x,y')
F(x,y) = vector([4*x*y,2*x^2])
show('int Fx in x = ', integrate(F(x,y)[0],x),
   '  int Fy in y = ', integrate(F(x,y)[1],y))

f(x,y) = 2*x^2*y
show('f(x,y) = ', f(x,y))
Int = f(1,1)-f(0,0)
show(' Int = ', Int)
\end{verbatim}
\end{programcode}
The output is shown in Fig \ref{PathIndepFdiff}
\begin{figure}[H]
\sidecaption[t]
\includegraphics[scale=.65]{chapters/ch10VectorFields/PathIndepFdiff}
\caption{Path independence, potential function}
\label{PathIndepFdiff}       % Give a unique label
\end{figure}

For example, for vector field $\mathbf{F}= \left< 3+2xy, x^2-3y^2 \right>$ let's find $\int_C  \mathbf{F}d\mathbf{r}$ for ANY path C from $(1,-2)$ to $(3,4)$.   First, let's perform partial integration to discover the potential function.
$$ \int M(x,y)dx= \int (3+2xy)dx = 3x+x^2y$$
$$  \int N(x,y)dx= \int (x^2-3y^2)dy = x^2y-y^3$$
Therefore, the potential function is $f(x,y) =x^2y-y^3+3x$.   The SageMath cell below illustrates this computation.

\begin{programcode}{SageMath Code and Output:}
\begin{verbatim}
reset(); var('x,y')
F(x,y) = [3+2*x*y, x^2 - 3*y^2]
show('int Fx in x = ', integrate(F(x,y)[0],x),
   '  int Fy in y = ', integrate(F(x,y)[1],y))

M(x,y) = F(x,y)[0];   N(x,y) = F(x,y)[1];
show('M(x,y) = ', M(x,y), '  N(x,y) = ', N(x,y))
show('dM/dy = ', M(x,y).diff(y),
     '    dN/dx = ', N(x,y).diff(x))

f(x,y) = x^2*y-y^3+3*x
Int = f(3,4)-f(1,-2)
show('f(x,y) = ',f(x,y),' Int = ', Int)
\end{verbatim}
\end{programcode}
The output is shown in Fig \ref{PathIndepFdiff2}
\begin{figure}[H]
\sidecaption[t]
\includegraphics[scale=.65]{chapters/ch10VectorFields/PathIndepFdiff2}
\caption{Path independence, potential function approach, no path given}
\label{PathIndepFdiff2}       % Give a unique label
\end{figure}

For a 3D example, let's come back to the conservative vector field $\mathbf{F} = \left< 2xy,x^2+2yz,y^2 \right>$ that we investigated before.  Just to remind
$$curl(\mathbf{F}) = \begin{pmatrix}
\mathbf{i}& \mathbf{j} & \mathbf{k}\\
\frac{\partial}{\partial x} & \frac{\partial}{\partial y} & \frac{\partial}{\partial z} \\
2xy & x^2+2yz & y^2
\end{pmatrix} =(2y-2y) \mathbf{i} - (0-0) \mathbf{j}+ (2x-2x) \mathbf{k} =   \mathbf{0} $$The partial integration shows:  $\int F_x dx = \int 2xy dx = x^2y $,  $\int F_y dy = \int x^2+2yz dy = x^2y+y^2z $, $\int F_z dz = \int y^2 dz = y^2z $, which  results in potential function $f(x,yz) =  x^2y+y^2z$.
Therefore, for any path from say $(1,-1,2)$ to $(4,5,7)$ the line integral $\int_L \mathbf{F} d \mathbf{r} = f(4,5,7) - f(1,-1,2) = 254$.

\begin{programcode}{SageMath Code and Output:}
\begin{verbatim}
reset(); var('x,y,z')
from sage.all import *
F(x,y,z) = vector([2*x*y,x^2+2*y*z,y^2])
show('F = ',F(x,y,z))
curlF(x,y,z) = F.curl()(x,y,z)
show('curl of F = ', curlF(x,y,z))
show('int Fx in x = ', integrate(F(x,y,z)[0],x))
show('int Fy in y = ', integrate(F(x,y,z)[1],y))
show('int Fz in z = ', integrate(F(x,y,z)[2],z))

f(x,y,z) = x^2*y+y^2*z
show('f(x,y,z) = ', f(x,y,z))
show(' grad of f = ',f.gradient()(x,y,z))
Int = f(4,5,7) - f(1,-1,2)
show('Int = ',Int)
\end{verbatim}
\end{programcode}
The output is shown in Fig \ref{PathIndepFdiff3D}
\begin{figure}[H]
\sidecaption[t]
\includegraphics[scale=.65]{chapters/ch10VectorFields/PathIndepFdiff3D}
\caption{Path independence, potential function approach in 3D}
\label{PathIndepFdiff3D}       % Give a unique label
\end{figure}

Finally, let's consider a case of a conservative vector field integral that can be computed much more easily on a different path.   Specifically, let  $\mathbf{F} = \left< M, N \right> =\left< 2xy^3+2, 3x^2y^2-3 \right>$, then the conservative field condition check confirms it: $\frac{\partial M}{\partial y} = 6xy^2 = \frac{\partial N}{\partial x} $.   The problem is to compute $\int_C \mathbf{F}d\mathbf{r}$ over the upper half circle C centered at $(1,0)$ of radius 1.   The direct path computation is illustrated in one of the SageMath cells below, but it is rather cumbersome with tedious high-power trigonometric integrals.   Two other approaches are much easier.   First of all, because it is a conservative vector field, we can use partial integration to find the potential function.  
$$\int M dx = x^2y^3+2x \ \ \ \  \int N dy = x^2y^3-3y \ \ \Rightarrow \ \ \ f(x,y) =  x^2y^3+2x-3y $$
Therefore,  $\int_C \mathbf{F}d\mathbf{r} = f(2,0)-f(0,0) = 4 - 0 = 4$.  Alternatively, using path independence, there is an even simpler way to compute along the straight line  $L: \ \mathbf{r}(t) = \left< t, 0 \right> \ \ 0\le t \le 2$.  Along the first path, $\mathbf{r}'(t) = \left< 1, 0\right>$, so $\mathbf{F}d\mathbf{r} = \left< 0+2,0-3 \right>\cdot\left< 1, 0 \right>=2$, therefore $\int_L \mathbf{F}d\mathbf{r} = \int_0^2 2 dt = 4$.   The SageMath cells below illustrate these computations.  First, a potential function approach:

\begin{programcode}{SageMath Code and Output:}
\begin{verbatim}
reset(); var('x,y')
F(x,y) = [2*x*y^3+2, 3*x^2*y^2-3]
show('int Fx in x = ', integrate(F(x,y)[0],x),
   '  int Fy in y = ', integrate(F(x,y)[1],y))

M(x,y) = F(x,y)[0];   N(x,y) = F(x,y)[1];
show('M(x,y) = ', M(x,y), '  N(x,y) = ', N(x,y))
show('dM/dy = ', M(x,y).diff(y),
     '    dN/dx = ', N(x,y).diff(x))

f(x,y) = x^2*y^3+2*x-3*y
show('f(x,y) = ', f(x,y))
show(' grad of f = ',f.gradient()(x,y))
Int = f(2,0) - f(0,0)
show('Int = ',Int)
\end{verbatim}
\end{programcode}
The output is shown in Fig \ref{PathIndep3}
\begin{figure}[H]
\sidecaption[t]
\includegraphics[scale=.65]{chapters/ch10VectorFields/PathIndep3}
\caption{Path independence, potential function approach in 3 ways}
\label{PathIndep3}       % Give a unique label
\end{figure}

Next, the direct computation along the circle:

\begin{programcode}{SageMath Code and Output:}
\begin{verbatim}
reset(); var('x,y')
F(x,y) = [2*x*y^3+2, 3*x^2*y^2-3]
show('int Fx in x = ', integrate(F(x,y)[0],x),
   '  int Fy in y = ', integrate(F(x,y)[1],y))

M(x,y) = F(x,y)[0];   N(x,y) = F(x,y)[1];
show('M(x,y) = ', M(x,y), '  N(x,y) = ', N(x,y))
show('dM/dy = ', M(x,y).diff(y),
     '    dN/dx = ', N(x,y).diff(x))

f(x,y) = x^2*y^3+2*x-3*y
show('f(x,y) = ', f(x,y))
show(' grad of f = ',f.gradient()(x,y))
Int = f(2,0) - f(0,0)
show('Int = ',Int)
\end{verbatim}
\end{programcode}
The output is shown in Fig \ref{PathIndepDirect}
\begin{figure}[H]
\sidecaption[t]
\includegraphics[scale=.65]{chapters/ch10VectorFields/PathIndepDirect}
\caption{Path independence, direct integration}
\label{PathIndepDirect}       % Give a unique label
\end{figure}

Finally, let's compute through the easiest horizontal line path.

\begin{programcode}{SageMath Code and Output:}
\begin{verbatim}
r(t) = vector([t,0])
t0=0; tf=2
rp(x,y,z) = derivative(r(t),t)

show('F(x,y,z) = ',F(x,y,z))
show('r(t) = ',r(t))
show('rp(t) = ',rp(t))

Fdr(t) = F(r(t)[0],r(t)[1]).dot_product(rp(t))
show('Fdr(t) = ',Fdr(t))
Int(t) = integrate( Fdr(t), t)
show('Int(t) = ', Int(t))
Int1 = Int(tf) - Int(t0)
show('Int1 = ',Int1)
p1 = plot_vector_field(F(x,y),(x,0,2),(y,0,1),
                       figsize = (4,3))
p2 = parametric_plot(r(t),(t,t0,tf),thickness=5)
(p1+p2).show()
\end{verbatim}
\end{programcode}
The output is shown in Fig \ref{PathIndep3Line}
\begin{figure}[H]
\sidecaption[t]
\includegraphics[scale=.65]{chapters/ch10VectorFields/PathIndep3Line}
\caption{Path independence, direct integration on horizontal line}
\label{PathIndep3Line}       % Give a unique label
\end{figure}










\section{Green's Theorem}

The Green's Theorem allows to compute line integrals over closed counterclockwise oriented path C bounding a simply connected region R as double integrals over R: $\int_C Mdx+Ndy = \int \int_R\left( \frac{\partial N}{\partial x} - \frac{\partial M}{\partial y} \right) dA$.  To illustrate it, let's compute both sides of Green's Theorem for  $M(x,y) = 3y^2$ and $N(x,y) = 2x^3+3xy$, then  $\frac{\partial N}{\partial x} - \frac{\partial M}{\partial y} = 6x^2+3y-6y= 6x^2-3y$.   Let the region R be bounded below by $L(x)=x^2$ and above by $U(x)=x$ for $0\le x\le 1$.   The line integral on the left has to be broken into two integrals: over $C_1:  \ y=x^2 \ \ x:0 \rightarrow 1$ and  $C_2: \ y=x \ \ x:1 \rightarrow 0$ in counterclockwise direction.  
$$ \int_{C_1} 3y^2dx+(2x^3+3xy)dy = \int_0^1  3(x^2)^2dx+(2x^3+3x x^2)2xdx = $$
$$ \int_0^1 13x^4 dx = \left( \frac{13x^5}{5} \right)_0^1=\frac{13}{5}$$
$$ \int_{C_2} 3y^2dx+(2x^3+3xy)dy = \int_1^0  3(x)^2dx+(2x^3+3x x)dx =$$
 $$ \int_1^0 2x^3+6x^2 dx =  \left(2 \frac{x^4}{4} + 6\frac{x^3}{3}\right)_1^0 = -\frac{5}{2}$$
Then the overall $\int_C =\int_{C_1} +  \int_{C_2} =\frac{1}{10}$.  

\begin{programcode}{SageMath Code and Output:}
\begin{verbatim}
reset(); var('x,y');
M(x,y) = 3*y^2;  N(x,y) = 2*x^3+3*x*y
show('M(x,y) = ', M(x,y), '  N(x,y) = ', N(x,y))

y1(x) = x^2;  x0=0; xf=1
y1p(x) = derivative(y1(x),x);
show('y1(x) = ',y1(x),'  y1p(x) = ',y1p(x))

Integrand = M(x,y1(x))+N(x,y1(x))*y1p(x)
show('Integrand = ',Integrand)
Int(x) = integrate( Integrand, x)
show('Int(x) = ', Int(x))
IntC1 = Int(xf) - Int(x0)
show('IntC1 = ',IntC1)
p1 = plot_vector_field([M(x,y),N(x,y)],(x,0,1),(y,0,1))
p2 = plot(y1(x),(x,x0,xf),figsize=(4,3))
(p1+p2).show()

##################################
y1(x) = x; y1p(x) = derivative(y1(x),x); x0=1; xf=0
show('y1(x) = ',y1(x),'  y1p(x) = ',y1p(x))

Integrand = M(x,y1(x))+N(x,y1(x))*y1p(x)
show('Integrand = ',Integrand)
Int(x) = integrate( Integrand, x)
show('Int(x) = ', Int(x))
IntC2 = Int(xf) - Int(x0)
show('IntC2 = ',IntC2)

show('IntC1 + IntC2 = ',IntC1+IntC2)
p1 = plot_vector_field([M(x,y),N(x,y)],(x,0,1),(y,0,1))
p2 = plot(y1(x),(x,0,1),figsize=(4,3))
(p1+p2).show()
\end{verbatim}
\end{programcode}
The output is shown in Fig \ref{GreenTheor1}
\begin{figure}[H]
\sidecaption[t]
\includegraphics[scale=.65]{chapters/ch10VectorFields/GreenTheor1}
\caption{Green's Theorem, line integrals computation}
\label{GreenTheor1}       % Give a unique label
\end{figure}

Let's now consider the right-hand side of the Green's Theorem.   It is a double integral:
$$ \int \int_R\left( \frac{\partial N}{\partial x} - \frac{\partial M}{\partial y} \right) dA = \int_0^1 \int_{x^2}^x 6x^2-3y dy dx =\int_0^1 \left( 6x^2y-3\frac{y^2}{2} \right)_{x^2}^xdx =  $$
$$ \int_0^1 6x^2(x-x^2)-3\frac{x^2-x^4}{2}dx   =  \int_0^1 -\frac{9}{2} \, x^{4} + 6 \, x^{3} - \frac{3}{2} \, x^{2}dx = $$
$$ \left( -\frac{9}{10} \, x^{5} + \frac{3}{2} \, x^{4} - \frac{1}{2} \, x^{3}  \right)_0^1 =\frac{1}{10} $$
This is exactly the same as the line integral over C, thus confirming Green's Theorem.

\begin{programcode}{SageMath Code and Output:}
\begin{verbatim}
U(x) = x # upper bound
L(x) = x^2 # lower bound
sol = solve(U(x)==L(x),x); 
show('Show intersection points: ',sol)
p = plot([L(x), U(x)], (x,0,1), 
     fill={0:[1]},figsize=(4,3)); show(p)

f(x,y) = derivative(N(x,y),x) - derivative(M(x,y),y) 
show('f(x,y) = ', f(x,y))
InnerInt(x,y) = integral(f(x,y),y)
show('InnerInt(x,y) = ', InnerInt(x,y))
Inner(x) = expand(InnerInt(x,U(x))-InnerInt(x,L(x)))
show('Inner(x)  = ',Inner(x) )
Int_outer(x) = integral(Inner(x),x)
Int = Int_outer(1) - Int_outer(0)
show(' Int_outer = ',Int_outer(x))
show(' Int = ',Int)

show('Sage Integral = ',
     f(x,y).integrate(y,L(x),U(x)).integrate(x,0,1))
\end{verbatim}
\end{programcode}
The output is shown in Fig \ref{GreenTheor2}
\begin{figure}[H]
\sidecaption[t]
\includegraphics[scale=.65]{chapters/ch10VectorFields/GreenTheor2}
\caption{Green's Theorem, area integral approach}
\label{GreenTheor2}       % Give a unique label
\end{figure}

As another example, let's consider line integral $\int_C (2xy + tan(\sqrt{x}))dx + (x^2+3x^2y-e^{-y^2})dy$ where path C is around 1st octant of parabola $y=4-x^2$ from $x=0$ to $x=2$.
The direct computation of three line integrals is not feasible. However, applying Green's Theorem 
$$\int_C Mdx+Ndy = \int_C (2xy + tan(\sqrt{x}))dx + (x^2+3x^2y-e^{-y^2})dy =$$
$$  \int \int_R\left( \frac{\partial N}{\partial x} - \frac{\partial M}{\partial y} \right) dA = \int_0^2 \int_0^{4-x^2} 6xy+2x-2x dy dx =\int_0^2  \left( 3xy^2\right)_0^{4-x^2}dx =  $$
$$ \int_0^2 3x(4-x^2)^2dx   =  \int_0^2 3x(16-8x^2+x^4)dx =  \int_0^2 48x-24x^3+3x^5dx = $$
$$ \left( 24x^2-6x^4+\frac{3}{6} x^6 \right)_0^2 =32 $$
Note that if the vector field is conservative, then $\frac{\partial N}{\partial x} - \frac{\partial M}{\partial y}=0$, so the line integral of a continuous conservative field over a closed path is 0.  For example, changing our example slightly we obtain  $\int_C (2xy + tan(\sqrt{x}))dx + (x^2-e^{-y^2})dy =  \int_0^2 \int_0^{4-x^2} 2x-2x dy dx=0$

\begin{programcode}{SageMath Code and Output:}
\begin{verbatim}
reset(); var('x,y');
M(x,y) = 2*x*y + tan(sqrt(x));  
N(x,y) = x^2+3*x^2*y-e^(-y^2)
show('M(x,y) = ', M(x,y), '  N(x,y) = ', N(x,y))
f(x,y) = derivative(N(x,y),x) - derivative(M(x,y),y) 
show('f(x,y) = ', f(x,y))
U(x) = 4-x^2 # upper bound
L(x) = 0 # lower bound
sol = solve(U(x)==L(x),x); 
show('Show intersection points: ',sol)
p = plot([L(x), U(x)], (x,0,2), 
    fill={0:[1]},figsize=(4,3)); show(p)

show('Sage Integral = ',
 f(x,y).integrate(y,L(x),U(x)).integrate(x,0,2))
\end{verbatim}
\end{programcode}
The output is shown in Fig \ref{GreenTheorArea}
\begin{figure}[H]
\sidecaption[t]
\includegraphics[scale=.65]{chapters/ch10VectorFields/GreenTheorArea}
\caption{Green's Theorem, area integral only possible}
\label{GreenTheorArea}       % Give a unique label
\end{figure}

Let's consider a similar kind of line integral over the upper half of an annulus with inner radius $1$ and upper radius $3$.
The Green's Theorem gives:
$$\int_C Mdx+Ndy = \int_C (2y^2 + \ln(x^2+1))dx + (x^3-\sin{y^2})dy = \int \int_R\left( \frac{\partial N}{\partial x} - \frac{\partial M}{\partial y} \right) dA =$$
$$   \int \int_R\left( 3x^2-4y \right) dA =  \int_0^{\pi} \int_1^3  (3r^2\cos^2(\theta) - 4r\sin(\theta) )r dr d\theta $$
$$    \int_0^{\pi}   \left(3\frac{r^4}{4}\cos^2(\theta) - 4\frac{r^3}{3}\sin(\theta) \right)_1^3 d\theta = \int_0^{\pi}   \left(60\cos^2(\theta) - \frac{104}{3}\sin(\theta) \right) d\theta $$
$$   \int_0^{\pi}   \left(60\frac{1+\cos(2\theta)}{2} - \frac{104}{3}\sin(\theta) \right) d\theta  = \left( 30t + 30\frac{\sin(2\theta)}{2} +  \frac{104}{3}\cos(\theta) \right)_0^{\pi} = 30\pi-\frac{208}{3}$$

\begin{programcode}{SageMath Code and Output:}
\begin{verbatim}
reset(); var('x,y,r,t');
M(x,y) = 2*y^2 + ln(x^2+1);  N(x,y) = x^3-sin(y^2)
show('M(x,y) = ', M(x,y), '  N(x,y) = ', N(x,y))
f(x,y) = derivative(N(x,y),x) - derivative(M(x,y),y) 
show('f(x,y) = ', f(x,y))

t = var('t')
p = polar_plot([1,3], (t, 0, pi), figsize = (4,3), 
           thickness=2, fill=True, fillcolor='orange')
show(p)

f(r,t) = ( 3*r^2*(cos(t))^2 - 4*r*sin(t) )*r
show('Sage Integral = ',
     f(r,t).integrate(r,1,3).integrate(t,0,pi))
\end{verbatim}
\end{programcode}
The output is shown in Fig \ref{GreenTheorAnnulus}
\begin{figure}[H]
\sidecaption[t]
\includegraphics[scale=.65]{chapters/ch10VectorFields/GreenTheorAnnulus}
\caption{Green's Theorem, over annular region}
\label{GreenTheorAnnulus}       % Give a unique label
\end{figure}

Green's Theorem can also be used in opposite directions to compute an integral line to find the area.   If we choose $M$ and $N$ functions such as $ \frac{\partial N}{\partial x} - \frac{\partial M}{\partial y} = 1$, then  $\int_C Mdx+Ndy = \int \int_R\left( \frac{\partial N}{\partial x} - \frac{\partial M}{\partial y} \right) dA =  \int \int_R\left(1 \right) dA  = A$.  One such choice is $M = -\frac{y}{2}$ and $N=\frac{x}{2}$.  For example, let's use a line integral to find the area of an ellipse.  We parametrize ellipse as $x=a\cos(t) \ \ \ y=b\sin(t) \ \ \ 0\le t \le 2\pi$.  Then:
$$A = \frac{1}{2} \int x dy-y dx=  \frac{1}{2}\int_0^{2\pi} a\cos(t)b\cos(t)dt-b\sin(t)(-a\sin(t)) dt =$$
$$ \frac{1}{2} \int_0^{2\pi} ab dt = \frac{1}{2}2\pi a b =\pi a b$$
The computation is a lot easier than direct trigonometric substitution integral.   The computation is rather simple, but we computed it in the SageMath cell below anyway.

\begin{programcode}{SageMath Code and Output:}
\begin{verbatim}
reset(); var('x,y,a,b,t')
M(x,y) = -y/2;  N(x,y) = x/2;  t0=0; tf=2*pi
x1(t) = a*cos(t);  y1(t) = b*sin(t)
x1p(t) = derivative(x1(t),t); y1p(t) = derivative(y1(t),t)

show('M(x,y) = ',M(x,y),'  N(x,y) = ',N(x,y))
show('x1(t) = ',x1(t),'  y1(t) = ',y1(t))
show('x1p(t) = ',x1p(t),'  y1p(t) = ',y1p(t))

Fdr(t) = (M(x1(t),y1(t))*x1p(t) + 
          N(x1(t),y1(t))*y1p(t) ).full_simplify()
show('Fdr(t) = ',Fdr(t))
Int(t) = integrate( Fdr(t), t)
show('Int(t) = ', Int(t))
Int1 = Int(tf) - Int(t0)
show('Int1 = ',Int1)
p1 = plot_vector_field([M(x,y),N(x,y)], 
         (x,-3,3), (y,-2,2), figsize=(4,3))
p2 = parametric_plot([3*cos(t),2*sin(t)],(t,0,2*pi))
(p1+p2).show()
\end{verbatim}
\end{programcode}
The output is shown in Fig \ref{GreenEllipse}
\begin{figure}[H]
\sidecaption[t]
\includegraphics[scale=.65]{chapters/ch10VectorFields/GreenEllipse}
\caption{Green's Theorem, ellipse integral}
\label{GreenEllipse}       % Give a unique label
\end{figure}

An interesting example is provided by vector field $\mathbf{F} =\left< M,N \right>=
 \left< -\frac{y}{x^2+y^2},\frac{x}{x^2+y^2}\right>$. Then:
$$\frac{\partial M}{\partial y} = -\frac{x^2+y^2 - y2y}{(x^2+y^2)^2} =\frac{y^2-x^2}{(x^2+y^2)^2} $$
$$\frac{\partial N}{\partial x} = \frac{x^2+y^2 - x2x}{(x^2+y^2)^2} =\frac{y^2-x^2}{(x^2+y^2)^2} $$
Therefore, any line integral over a closed path C is 0:
$$\int_C Mdx + Ndy = \int\int_R \frac{\partial N}{\partial x} - \frac{\partial M}{\partial y}dA = 0$$
Consider, however, a circle of arbitrary radius $a$ centered at the origin.  Parametrize it as $x=a\cos(t) \ \ \
y = b\sin(t) \ \ \ 0\le t\le 2\pi$, so $x^2+y^2=a^2$.   Then
$$\int_C Mdx + Ndy = \int_0^{2\pi} \left( \frac{-a\sin(t)}{a^2}a(-\sin(t)) + \frac{a\cos(t)}{a^2}a\cos(t) \right)dt =  $$
$$\int_0^{2\pi} \sin^2(t) + \cos^2(t) dt = 2\pi \ne 0$$
There is no contradiction here because Green's Theorem stipulates that vector field $\mathbf{F}$ is continuous
in domain R bounded by closed curve C, but the field is undefined at the origin.  However, for any closed curve not
containing the origin, the closed loop integral is indeed 0.

Consider now any closed path $C_1$ enclosing the origin.
We can always find a counterclockwise traversed circle $C$ of a small enough radius to be completely inside of this curve.  Then the region $D$ between $C_1$ and $C$ is bounded by positively \(counterclockwise\) oriented $C_1$ curve and negatively \(clockwise\)
oriented small circle $C$.  One can think of a small cut in between them connected back and forth by straight lines.  The region $D$ does not contain the origin, so our vector field is continuous in it.  Therefore
$$\int_{C_1}Mdx + Ndy+\int_{-C}Mdx + Ndy =  \int\int_D \frac{\partial N}{\partial x} - \frac{\partial M}{\partial y}dA = 0$$
Thus, regardless of the complexity of the path $C_1$:   $$ \int_{C_1} Mdx + Ndy = -\int_{-C}Mdx + Ndy = \int_C Mdx + Ndy= 2\pi$$The SageMath cell below illustrates the vector field and circle integral computation.  We also indicate $C_1$ as an ellipse.

\begin{programcode}{SageMath Code and Output:}
\begin{verbatim}
reset(); var('x,y,a,t')
M(x,y) = -y/(x^2+y^2);  N(x,y) = x/(x^2+y^2);  t0=0; tf=2*pi
x1(t) = a*cos(t);  y1(t) = a*sin(t)
x1p(t) = derivative(x1(t),t); y1p(t) = derivative(y1(t),t)

show('M(x,y) = ',M(x,y),'  N(x,y) = ',N(x,y))
show('dM/dy = ', M(x,y).diff(y).full_simplify(),
     '    dN/dx = ', N(x,y).diff(x).full_simplify())
show('x1(t) = ',x1(t),'  y1(t) = ',y1(t))
show('x1p(t) = ',x1p(t),'  y1p(t) = ',y1p(t))

Fdr(t) = (M(x1(t),y1(t))*x1p(t) + 
          N(x1(t),y1(t))*y1p(t) ).full_simplify()
show('Fdr(t) = ',Fdr(t))
Int(t) = integrate( Fdr(t), t)  
show('Int(t) = ', Int(t))
Int1 = Int(tf) - Int(t0)
show('Int1 = ',Int1)
p1 = plot_vector_field([M(x,y),N(x,y)], 
           (x,-3,3), (y,-2,2),figsize=(4,3))
p2 = parametric_plot([cos(t),sin(t)],(t,0,2*pi))
p3 = parametric_plot([3*cos(t),2*sin(t)],(t,0,2*pi))
(p1+p2+p3).show()
\end{verbatim}
\end{programcode}
The output is shown in Fig \ref{GreenInfOrigin}
\begin{figure}[H]
\sidecaption[t]
\includegraphics[scale=.65]{chapters/ch10VectorFields/GreenInfOrigin}
\caption{Green's Theorem, field infinite at the origin}
\label{GreenInfOrigin}       % Give a unique label
\end{figure}









\section{Parametric Surfaces}

Let's begin with a basic parametric surface given by $x=4\cos(u)\ \ \ y=4\sin(u) \ \ \ z = v \ \ 0\le u \le 2\pi \ \ 0 \le v \le 7$.
Note that $x^2+y^2=16$ - circle of radius 4 centered at the origin, so $u$ plays the role of the polar angle $\theta$. On the
other hand, $z=v$, so $v$ is just the height.   This surface represents a cylinder along the z axis of radius $4$ stretching
from $0$ to $7$. A simple switch of variables would produce a cylinder along some other coordinate axis. For example:
$x=4\cos(u)\ \ \ y=v \ \ \ z = 4\sin(u) \ \ 0\le u \le 2\pi \ \ 0 \le v \le 7$ is the same cylinder along $y$ axis.
The SageMath cell below shows these cylinders with \texttt{parametric\_plot3d}, which was also used for 3D curves.

\begin{programcode}{SageMath Code and Output:}
\begin{verbatim}
var('u,v')
p1 = parametric_plot3d([4*cos(u),4*sin(u),v],
           (u,0,2*pi),(v,0,7),opacity=0.5)
p2 = parametric_plot3d([4*cos(u),v,4*sin(u)],
           (u,0,2*pi),(v,0,7),opacity=0.5,color='red')
(p1+p2).show()
\end{verbatim}
\end{programcode}
The output is shown in Fig \ref{TwoCylinders}
\begin{figure}[H]
\sidecaption[t]
\includegraphics[scale=.65]{chapters/ch10VectorFields/TwoCylinders}
\caption{Two cylinders interseting}
\label{TwoCylinders}       % Give a unique label
\end{figure}

Another standard parametric representation is $x=4\sin(u)\cos(v)\ \ \ y=4\sin(u)\sin(v) \ \ \ z = 4\cos(u) \ \ 0\le u
\le \pi \ \ 0 \le v \le 2\pi$.  It is just a sphere of radius $4$ centered in the origin written in terms of
spherical coordinates $u=\phi$ and $v=\theta$, while $\rho=4=const$.  
$$x^2+y^2+z^2=16\sin^2(u)\cos^2(v)+16\sin^2(u)\sin^2(v)+16\cos^2(u)=$$
$$16\sin^2(u)+16\cos^2(u)=16$$In the SageMath cell below we illustrate this sphere as well as a portion of a smaller radius sphere in the first quadrant.

\begin{programcode}{SageMath Code and Output:}
\begin{verbatim}
restart; var('u,v')
p1 = parametric_plot3d(
    [4*sin(u)*cos(v),4*sin(u)*sin(v),4*cos(u)],
     (u,0,pi),(v,0,2*pi),opacity=0.5)
p2 = parametric_plot3d(
    [3*sin(u)*cos(v),3*sin(u)*sin(v),3*cos(u)],
 (u,0,pi/2),(v,0,pi/2),color='red',opacity=0.5)
(p1+p2).show()
\end{verbatim}
\end{programcode}
The output is shown in Fig \ref{ParametricSphere}
\begin{figure}[H]
\sidecaption[t]
\includegraphics[scale=.65]{chapters/ch10VectorFields/ParametricSphere}
\caption{Parametric representation of sphere}
\label{ParametricSphere}       % Give a unique label
\end{figure}

If the explicit surface equation is available, a parametric representation is straightforward.  For example,
for a standard cone along $z$-axis: $z=\sqrt{x^2+y^2}$, the direct parametric representation is $x=u; \ \ y = v;
 \ \ z = \sqrt{u^2+v^2}$.  By far, a parametric representation is not unique. Say, we can use cylindrical
coordinates to rewrite the cone as $x = u\cos(v); \ \ y = u\sin(v); \ \ z = u$, where $u=r$ and $v=\theta$. Analogously,
for a paraboloid $z=x^2+y^2$, the direct parametrization is $x=u; \ \ y = v;  \ \ z = u^2+v^2$, while the cylindrical
parametrization is $x = u\cos(v); \ \ y = u\sin(v); \ \ z = u^2$.  Both surfaces are illustrated in the SageMath cell
below.

\begin{programcode}{SageMath Code and Output:}
\begin{verbatim}
restart; var('u,v')
p1 = parametric_plot3d([u,v,sqrt(u^2+v^2)],
            (u,-2,2),(v,-2,2),opacity=0.5)
p2 = parametric_plot3d([u,v,u^2+v^2],
 (u,-2,2),(v,-2,2),color='red',opacity=0.5)
(p1+p2).show()
\end{verbatim}
\end{programcode}
The output is shown in Fig \ref{ParametricExplicit}
\begin{figure}[H]
\sidecaption[t]
\includegraphics[scale=.65]{chapters/ch10VectorFields/ParametricExplicit}
\caption{Parametric explicit surface equations}
\label{ParametricExplicit}       % Give a unique label
\end{figure}

Let's now compute normal vector and tangent line for the paraboloid above $\mathbf{r}(u,v) = \left< u,v,u^2+v^2 \right>$.  The partial derivative tangent vectors are $\mathbf{r}_u = \left< 1,0,2u \right>$ and $\mathbf{r}_v = \left< 0,1,2v \right>$, therefore normal vector is:
$$ \mathbf{n} = \mathbf{r}_u \times \mathbf{r}_v = \begin{pmatrix}
\mathbf{i} & \mathbf{j} & \mathbf{k}\\
1 & 0 & 2u \\
0 & 1 & 2v
\end{pmatrix} =(-2u) \mathbf{i}-(2v) \mathbf{j} +(1-0) \mathbf{k} = \left<  -2u,-2v,1 \right>$$
At a particular point corresponding to $u=2; \ v = 1$, we obtain $\mathbf{r}(2,1) = \left< 2,1,5 \right>$ and normal vector $\mathbf{n} = \left< -4,-2,1 \right> $.  Therefore, the tangent plane is:
$$-4(x-2)-2(y-1)+1(z-5)=0 \ \ \Rightarrow \ \ $$
$$ z = 5+4(x-2)+2(y-1) = 4x+2y-5$$
The normal line is $x=2-4t; \ \ y=1-2t; \ \ z = 5+t$.

\begin{programcode}{SageMath Code and Output:}
\begin{verbatim}
reset; var('u,v,t,x,y,z'); 
r(u,v) = vector([u,v,u^2+v^2])
ru(u,v) = derivative(r(u,v),u)
rv(u,v) = derivative(r(u,v),v)
show('r(u,v) = ',r(u,v))
show(' ru(u,v) = ',ru(u,v),' rv(u,v) = ',rv(u,v))
n(u,v) = ru(u,v).cross_product(rv(u,v))
show('normal vector n(u,v) = ',n(u,v))

u1 = 2; v1 = 1; x1 = u1; y1 = v1; z1 = u1^2+v1^2
ru1 = ru(u1,v1); rv1 = rv(u1,v1); n1 = n(u1,v1)
show('u1 = ',u1,' v1 = ',v1,' x1 = ',
           x1,' y1 = ',y1,' z1 = ',z1)
show('ru1 = ',ru1,' rv1 = ',rv1,' n1 = ',n1)

tang_plane = n1[0]*(x-x1)+n1[1]*(y-y1)+ n1[2]*(z-z1);
normal_line = [x1+n1[0]*t,y1+n1[1]*t,z1+n1[2]*t]
tang_plane_angle = arccos(abs(n1[2])/n1.norm(2))
show('tang_plane = ', tang_plane)
show('normal_line = ', normal_line)
show('tangent plane angle = ',N(tang_plane_angle*180/pi))
p1 = parametric_plot3d(r(u,v),(u,-3,3),(v,-3,3),opacity=0.5)
p2 = implicit_plot3d(tang_plane,(x,-3,3),(y,-3,3),(z,0,7),
                    color = 'green',opacity = 0.5)
p3 = parametric_plot3d(normal_line,(t,-0.5,0.5),
                       color='red',thickness=5)
(p1+p2+p3).show()
\end{verbatim}
\end{programcode}
The output is shown in Fig \ref{Param3DTangNorm}
\begin{figure}[H]
\sidecaption[t]
\includegraphics[scale=.65]{chapters/ch10VectorFields/Param3DTangNorm}
\caption{Parametric paraboloid tangent plane and normal line}
\label{Param3DTangNorm}       % Give a unique label
\end{figure}

The general formula for the surface area of the parametric surface is given by $S = \int\int || \mathbf{r}_u \times \mathbf{r}_v || dvdu$.  This often involves rather complex integrals unless there is an underlying symmetry.  For example in the case of a paraboloid above, let's find the area of the paraboloid bounded by $z=4$ plane.   It results in $z=u^2+v^2=4$, so polar coordinates can be used in the integration.  In addition, the integrand simplifies to $|| \mathbf{r}_u \times \mathbf{r}_v || = \sqrt{4u^2+4v^2+1} =\sqrt{4r^2+1}$, therefore
$$S = \int\int || \mathbf{r}_u \times \mathbf{r}_v || dvdu = \int_0^{2\pi}\int_0^2 \sqrt{4r^2+1}r dr d\theta = [u=4r^2+1 \ \Rightarrow \ du=8rdr] = $$
$$ 2\pi \int_1^{17}\sqrt{u}\frac{du}{8} = \frac{\pi}{4} \left( \frac{u^{3/2}}{3/2} \right)_1^{17}=\frac{1}{6} \, \pi {\left(17 \, \sqrt{17} - 1\right)}$$
Note that in this case, $z=f(x,y)$ is given explicitly and we can write $\mathbf{r}(x,y) = \left< x,y,f(x,y) \right>$, so $\mathbf{r}_x(x,y) = \left< 1,0,f_x(x,y) \right>$ and $\mathbf{r}_y(x,y) = \left< 0,1,f_y(x,y) \right>$.  Then 
$$ \mathbf{n} = \mathbf{r}_x \times \mathbf{r}_y = \begin{pmatrix}
\mathbf{i} & \mathbf{j} & \mathbf{k}\\
1 & 0 & f_x(x,y) \\
0 & 1 & f_y(x,y)
\end{pmatrix} =(- f_x(x,y)) \mathbf{i}-( f_y(x,y)) \mathbf{j} +(1-0) \mathbf{k} = $$
$$ \left<  -f_x(x,y),-f_y(x,y),1 \right> \ \ \Rightarrow ||\mathbf{n}|| = ||\mathbf{r}_x \times \mathbf{r}_y|| = \sqrt{f_x(x,y)^2+f_y(x,y)^2+1}$$
This results in the standard surface area formula $S = \int\int  \sqrt{f_x(x,y)^2+f_y(x,y)^2+1} \ dA$

\begin{programcode}{SageMath Code and Output:}
\begin{verbatim}
var('r,t')
nnorm(r,t) = sqrt(1+4*r^2)
S = ( nnorm(r,t)*r ).integrate(r,0,2).integrate(t,0,2*pi)
show('nnorm(r,t) = ',nnorm(r,t),'  S = ',S)
\end{verbatim}
\end{programcode}
The output is shown in Fig \ref{ParamSurfaceArea}
\begin{figure}[H]
\sidecaption[t]
\includegraphics[scale=.65]{chapters/ch10VectorFields/ParamSurfaceArea}
\caption{Surface area of a parametric paraboloid}
\label{ParamSurfaceArea}       % Give a unique label
\end{figure}

Let's consider a more complex torus parametric surface: 
$$\mathbf{r}(u,v)  = \left< (5+2cos(u))cos(v),(5+2cos(u))sin(v),2sin(u) \right> $$
The calculations are more involved, but the SageMath cell below handles them as well.

\begin{programcode}{SageMath Code and Output:}
\begin{verbatim}
reset; var('u,v,t,x,y,z'); 
r(u,v) = vector([(5+2*cos(u))*cos(v),
        (5+2*cos(u))*sin(v),2*sin(u)])
ru(u,v) = derivative(r(u,v),u)
rv(u,v) = derivative(r(u,v),v)
show('r(u,v) = ',r(u,v))
show('ru(u,v) = ',ru(u,v))
show('rv(u,v) = ',rv(u,v))
n(u,v) = ru(u,v).cross_product(rv(u,v)) 
show('normal vector n(u,v) = <')
show(n(u,v)[0],','); show(n(u,v)[1],','); 
show(n(u,v)[2],'>')

u1 = pi/3; v1 = pi/4
x1 = (5+2*cos(u1))*cos(v1); 
y1 = (5+2*cos(u1))*sin(v1);
z1 = 2*sin(u1); ru1 = ru(u1,v1); 
rv1 = rv(u1,v1); n1 = n(u1,v1)
show('u1 = ',u1,' v1 = ',v1,' x1 = ',
           x1,' y1 = ',y1,' z1 = ',z1)
show('ru1 = ',ru1,' rv1 = ',rv1)
show(' n1 = ',n1)

tang_plane = n1[0]*(x-x1)+n1[1]*(y-y1)+ n1[2]*(z-z1);
normal_line = [x1+n1[0]*t,y1+n1[1]*t,z1+n1[2]*t]
tang_plane_angle = arccos(abs(n1[2])/n1.norm(2))
show('tang_plane = ', tang_plane)
show('normal_line = ', normal_line)
show('tangent plane angle = ',
     N(tang_plane_angle*180/pi))

#### Surface area computation
nnorm(u,v) = (n.norm(2)).full_simplify()
S = (nnorm(u,v)).integrate(u,0,2*pi).integrate(v,0,2*pi)
show('nnorm(u,v) = ',nnorm(u,v),'  S = ',S)

p1 = parametric_plot3d(r(u,v),(u,0,2*pi),
                       (v,0,2*pi),opacity=0.5)
p2 = implicit_plot3d(tang_plane,(x,-5,5),
   (y,-5,5),(z,-5,5),color = 'green',opacity = 0.5)
p3 = parametric_plot3d(normal_line,(t,-0.3,0.3),
                       color='red',thickness=5)
(p1+p2+p3).show()
\end{verbatim}
\end{programcode}
The output is shown in Fig \ref{TorusParametric}
\begin{figure}[H]
\sidecaption[t]
\includegraphics[scale=.65]{chapters/ch10VectorFields/TorusParametric}
\caption{Torus tangent plane, normal line, and surface area}
\label{TorusParametric}       % Give a unique label
\end{figure}

Lastly, let's consider a surface of rotation.   Assume a curve $y=f(x)$ is rotated about $x$\-axis.  We can represent it as $\mathbf{r}(u,v) = \left< u, f(u)cos(v), f(u)sin(v) \right>$.  
Then, partial derivative tangent vectors are $\mathbf{r}_u = \left< 1, f'(u)cos(v), f'(u)sin(v)  \right>$ and $\mathbf{r}_v = \left< 0, -f(u)sin(v), f(u)cos(v) \right>$, therefore normal vector is:
$$ \mathbf{n} = \mathbf{r}_u \times \mathbf{r}_v = \begin{pmatrix}
\mathbf{i} & \mathbf{j} & \mathbf{k}\\
1 & f'(u)cos(v) & f'(u)sin(v) \\
0 & -f(u)sin(v) & f(u)cos(v)
\end{pmatrix} = $$
$$[f(u)f'(u)(cos^2(v)+sin^2(v))] \mathbf{i}-(f(u)cos(v)) \mathbf{j} +(-f(u)sin(v)) \mathbf{k} =$$
$$ \left<  f(u)f'(u),-f(u)cos(v),-f(u)sin(v) \right>$$
Then the norm of this normal vector is:
$$||\mathbf{n}|| = \sqrt{ (f(u))^2(f'(v))^2+ (f(u))^2cos^2(v)+ (f(u))^2sin^2(v) } = f(u) \sqrt{ (f'(u))^2+ 1} $$
Therefore, the surface integral is
$$S =  \int\int || \mathbf{n} || dvdu =  \int_a^b\int_0^{2\pi} f(u) \sqrt{ (f'(u))^2+ 1} dvdu=  2\pi \int_a^b f(x) \sqrt{ (f'(x))^2+ 1} dx$$
For example, let $f(x)=\frac{1}{x}$, then $f'(x)=-\frac{1}{x^2}$.  
$$ S = 2\pi \int_0^1\frac{1}{x}\sqrt{ \frac{1}{x^4}+ 1}dx  $$The SageMath cell below plots the surface of rotation as well as tangent plane and normal line at $\left(u=1,v=\frac{\pi}{2}\right)$ and computes surface area of rotated segment from $x=1$ to $x=2$.

\begin{programcode}{SageMath Code and Output:}
\begin{verbatim}
reset; var('u,v,t,x,y,z'); assume(u>0)
r(u,v) = vector([u,1/u*cos(v),1/u*sin(v)])
ru(u,v) = derivative(r(u,v),u)
rv(u,v) = derivative(r(u,v),v)
show('r(u,v) = ',r(u,v))
show('ru(u,v) = ',ru(u,v))
show('rv(u,v) = ',rv(u,v))
n(u,v) = ru(u,v).cross_product(rv(u,v)) 
show('normal vector n(u,v) = ',n(u,v))

u1 = 1; v1 = pi/2
x1 = u1; y1 = 1/u1*cos(v1); z1 = 1/u1*sin(v1);
ru1 = ru(u1,v1); rv1 = rv(u1,v1); n1 = n(u1,v1)
show('u1 = ',u1,' v1 = ',v1,' x1 = ',
           x1,' y1 = ',y1,' z1 = ',z1)
show('ru1 = ',ru1,' rv1 = ',rv1)
show(' n1 = ',n1)

tang_plane = n1[0]*(x-x1)+n1[1]*(y-y1)+ n1[2]*(z-z1);
normal_line = [x1+n1[0]*t,y1+n1[1]*t,z1+n1[2]*t]
tang_plane_angle = arccos(abs(n1[2])/n1.norm(2))
show('tang_plane = ', tang_plane)
show('normal_line = ', normal_line)
show('tangent plane angle = ',N(tang_plane_angle*180/pi))

show('###### Surface area computation')
nnorm(u,v) = (n.norm(2)).full_simplify()
show((nnorm(u,v)).integrate(u).integrate(v))
S = (nnorm(u,v)).integrate(u,1,2).integrate(v,0,2*pi)
S = N(S)
show('nnorm(u,v) = ',nnorm(u,v),'  S = ',S)

p1 = parametric_plot3d(r(u,v),(u,0.5,3),(v,0,2*pi),opacity=0.5)
p2 = implicit_plot3d(tang_plane,(x,-2,2),(y,-2,2),(z,0,2),
                    color = 'green',opacity = 0.5)
p3 = parametric_plot3d(normal_line,(t,-0.3,0.3),
                       color='red',thickness=5)
(p1+p2+p3).show()
\end{verbatim}
\end{programcode}
The output is shown in Fig \ref{RotationSurface}
\begin{figure}[H]
\sidecaption[t]
\includegraphics[scale=.65]{chapters/ch10VectorFields/RotationSurface}
\caption{Surface of rotation tangent plane, normal line, and surface area}
\label{RotationSurface}       % Give a unique label
\end{figure}









\section{Surface Integrals}

Let's begin this section by evaluating the surface integral of $f(x,y) = x^2z-3y^2$ on the first octant portion of the plane $x + 2y + z = 2$.   In this case, it does not matter much, but it is more natural to project onto $xy$-plane.  Plugging $z=0$, we obtain $x + 2y = 2 \ \Rightarrow \ y=\frac{2-x}{2}$.  The surface equation is $z = g(x,y)=2-x-2y$, so $g_x(x,y)=1$  and $g_y(x,y)=2$.  Then the surface element square root term is just a number: $\sqrt{1+g_x^2+g_y^2} = \sqrt{1+1+4} =\sqrt{6}$.

Therefore 
$$\int\int_S f dS = \int\int_S f(x,y,g(x,y)) \sqrt{1+g_x^2+g_y^2}dA =  $$
$$\int_0^2 \int_0^{(2-x)/2} (x^2(2-x-2y)-3y^2) \sqrt{6} dy dx = -\frac{7\sqrt{6}}{30} $$

\begin{programcode}{SageMath Code and Output:}
\begin{verbatim}
reset(); var('x,y,z');
f(x,y,z) = x^2*z-3*y^2;  g(x,y) = 2-x-2*y;
show( implicit_plot(g(x,y)==0,(x,0,3),\
        (y,0,2), gridlines=True, figsize = (4,3)))

gx(x,y) = derivative(g(x,y),x)
gy(x,y) = derivative(g(x,y),y)
dS(x,y) = sqrt(1+gx(x,y)^2+gy(x,y)^2)
show('gx(x,y) = ',gx(x,y),' gy(x,y) = ',
     gy(x,y), ' dS = ', dS(x,y))
show('integrand = f(x,y,g(x,y))*dS = ', 
               f(x,y,g(x,y))*dS(x,y))
Int = ( f(x,y,g(x,y))*dS(x,y) ). \
        integrate(y,0,(2-x)/2).integrate(x,0,2)
show('Int = ', Int)
\end{verbatim}
\end{programcode}
The output is shown in Figure \ref{SurfIntPlane}
\begin{figure}[H]
\sidecaption[t]
\includegraphics[scale=.65]{chapters/ch10VectorFields/SurfIntPlane}
\caption{Surface integral on a tilted plane}
\label{SurfIntPlane}       % Give a unique label
\end{figure}


For some problems, we have to project to a different coordinate plane.  For example, find the surface integral of $f(x,y) = 2x+3y+4z$ on the first-octant part of a circular cylinder $x=g(y,z)=\sqrt{9-z^2}$ between $y=0$ and $y=4$.   In this case, we have to project onto $ yz$-plane, and the projection is just a rectangle. Plugging $x=0$, we obtain $\sqrt{9-z^2} = 0 \ \Rightarrow \ z=3$.  The surface equations derivatives are $g_y(y,z)=0$  and $g_z(y,z)=-\frac{z}{\sqrt{9-z^2}}$.  Then the surface element is: $\sqrt{1+g_y^2+g_z^2} = \sqrt{1+0+\frac{z^2}{9-z^2}} =\frac{3}{\sqrt{9-z^2}}$.   Therefore 
$$\int\int_S f dS = \int\int_S f(g(y,z),y,z) \sqrt{1+g_y^2+g_z^2}dA =  $$
$$\int_0^4 \int_0^3 (2\sqrt{9-z^2}+3y+4z) \frac{3}{\sqrt{9-z^2}} dz dy  $$
The integrals are somewhat tedious, but SageMath computes them.

\begin{programcode}{SageMath Code and Output:}
\begin{verbatim}
reset(); var('x,y,z');  assume(z<3)
f(x,y,z) =2*x+3*y+4*z;  g(y,z) = sqrt(9-z^2);
gy(y,z) = derivative(g(y,z),y)
gz(y,z) = derivative(g(y,z),z)
dS(y,z) = sqrt(1+gy(y,z)^2+gz(y,z)^2).full_simplify()
show('gy(y,z) = ',gy(y,z),' gz(y,z) = ',
     gz(y,z), ' dS = ', dS(y,z))
show('integrand = f(g(y,z),y,z)*dS = ', 
              f(g(y,z),y,z)*dS(y,z))
Int = ( f(g(y,z),y,z)*dS(y,z) ). \
        integrate(z,0,3).integrate(y,0,4)
show('Int = ', Int, ' = ', N(Int))
\end{verbatim}
\end{programcode}
The output is shown in Figure \ref{CylinderYZ}
\begin{figure}[H]
\sidecaption[t]
\includegraphics[scale=.65]{chapters/ch10VectorFields/CylinderYZ}
\caption{Surface integral in y-z plane}
\label{CylinderYZ}       % Give a unique label
\end{figure}


Actually, it is much easier to compute above surface integral using the parametrization  $\mathbf{r}(u,v) = \left< 3\cos(v),u,3\sin(v) \right> \ \ 0 \le u \le 4, \ \ 0 \le v \le \frac{\pi}{2}$.   Here $x^2+z^2 = 9\cos^2(v)+9\sin^2(v)=9$.
Then, partial derivative tangent vectors are $\mathbf{r}_u = \left< 0, 1,0  \right>$ and $\mathbf{r}_v = \left< -3\sin(v), 0, 3\cos(v) \right>$, therefore normal vector is:
$$ \mathbf{n} = \mathbf{r}_u \times \mathbf{r}_v = \begin{pmatrix}
\mathbf{i} & \mathbf{j} & \mathbf{k}\\
0           & 1 &     0 \\
-3\sin(v) & 0 & 3\cos(v)
\end{pmatrix} = $$
$$3\cos(v) \mathbf{i}-0 \mathbf{j} +(3\sin(v)) \mathbf{k} = \left<  3\cos(v),0,3\sin(v) \right>$$
Then the norm of this normal vector is:
$$||\mathbf{n}|| = \sqrt{9 \cos^2(v)+ 9\sin^2(v) } =3$$
Therefore, the surface integral is
$$\int\int_S fdS =  \int\int f || \mathbf{n} || dvdu =  \int_0^4\int_0^{\pi/2} [2(3\cos(v))+3u+4(3\sin(v))]3dvdu=$$
$$ \int_0^{\pi/2} \left(18u\cos(v)+9\frac{u^2}{2}+ 36u\sin(v))\right)_0^4  dv= $$
$$ \left(72\sin(v)+72v - 144\cos(v) \right)_0^{\pi/2} = 36\pi+216$$


\begin{programcode}{SageMath Code and Output:}
\begin{verbatim}
reset(); var('u,v,x,y,z'); assume(u>0)
f(x,y,z) =2*x+3*y+4*z;
r(u,v) = vector([3*cos(v),u,3*sin(v)])
ru(u,v) = derivative(r(u,v),u)
rv(u,v) = derivative(r(u,v),v)
show('r(u,v) = ',r(u,v))
show('ru(u,v) = ',ru(u,v))
show('rv(u,v) = ',rv(u,v))
n(u,v) = ru(u,v).cross_product(rv(u,v)) 
show('normal vector n(u,v) = ',n(u,v))
nnorm(u,v) = (n.norm(2)).full_simplify()
show('normal vector nnorm(u,v) = ',nnorm(u,v))
show(f(r(u,v)[0],r(u,v)[1],r(u,v)[2])* nnorm(u,v))
Int = ( f(r(u,v)[0],r(u,v)[1],r(u,v)[2])* nnorm(u,v) ). \
        integrate(u,0,4).integrate(v,0,pi/2)
show('Int = ', Int, ' = ', N(Int))
parametric_plot3d(r(u,v),(u,0,4),(v,0,pi/2),opacity=0.5)
\end{verbatim}
\end{programcode}
The output is shown in Figure \ref{CylinderYZ2}
\begin{figure}[H]
\sidecaption[t]
\includegraphics[scale=.65]{chapters/ch10VectorFields/CylinderYZ2}
\caption{Cylindrical parametrization of the above integral}
\label{CylinderYZ2}       % Give a unique label
\end{figure}



Let's now consider surface integral of $f(x,y,z)=x^2$ over a sphere of radius $a$ parametrized with spherical coordinates: $\mathbf{r}(u,v) = \left< a sin(u)cos(v),a sin(u)sin(v),a cos(u) \right> \ \ 0 \le u \le \pi, \ \ 0 \le v \le 2\pi$.  
Then, partial derivative tangent vectors are $\mathbf{r}_u = \left< a cos(u)cos(v),a cos(u)sin(v),-a sin(u)   \right>$ and $\mathbf{r}_v = \left<  -a sin(u)sin(v),a sin(u)cos(v),0 \right>$, therefore normal vector is:
$$ \mathbf{n} = \mathbf{r}_u \times \mathbf{r}_v = \begin{pmatrix}
\mathbf{i} & \mathbf{j} & \mathbf{k}\\
a cos(u)cos(v) & a cos(u)sin(v) & -a sin(u) \\
-a sin(u)sin(v) & a sin(u)cos(v)  &   0
\end{pmatrix} = $$
$$\left< a^2 sin^2(u)cos(v), a^2 sin^2(u)sin(v), a^2 cos(u)sin(u)(cos^2(v)+sin^2(v)) \right> $$
Then the norm of this normal vector is:
$$||\mathbf{n}|| = a^2\sqrt{sin^4(u)(cos^2(v)+sin^2(v))+cos^2(u)sin^2(u) } =a^2sin(u)$$
Therefore, surface integral is
$$\int\int_S fdS =  \int\int f || \mathbf{n} || dvdu =  \int_0^{2\pi}\int_0^{\pi} a^2 sin^2(u)cos^2(v) a^2sin(u)dvdu=$$
$$a^4 \int_0^{2\pi}cos^2(v) dv \int_0^{\pi} sin^3(u)du = a^4 \int_0^{2\pi} \frac{1+cos(2v)}{2} dv \int_0^{\pi} (cos^2(u)-1)d(cos(u)) =$$
$$ a^4 \left( \frac{u}{2}+\frac{sin(2v)}{4} \right)_0^{2\pi}  \left( \frac{cos^3(u)}{3} - cos(u) \right)_0^{\pi} =  \frac{4}{3}\pi a^4$$


\begin{programcode}{SageMath Code and Output:}
\begin{verbatim}
reset(); var('u,v,x,y,z,a'); assume(a>0); assume(sin(u)>0)
f(x,y,z) =x^2;
r(u,v) = vector([a*sin(u)*cos(v),a*sin(u)*sin(v),a*cos(u)])
ru(u,v) = derivative(r(u,v),u)
rv(u,v) = derivative(r(u,v),v)
show('r(u,v) = ',r(u,v))
show('ru(u,v) = ',ru(u,v))
show('rv(u,v) = ',rv(u,v))
n(u,v) = ru(u,v).cross_product(rv(u,v)) 
show('normal vector n(u,v) = <')
show(n(u,v)[0],','); show(n(u,v)[1],','); show(n(u,v)[2],'>')
nnorm(u,v) = (n.norm(2)).full_simplify()
show('normal vector nnorm(u,v) = ',nnorm(u,v))
show(f(r(u,v)[0],r(u,v)[1],r(u,v)[2])* nnorm(u,v))
Int = ( f(r(u,v)[0],r(u,v)[1],r(u,v)[2])* nnorm(u,v) ). \
        integrate(u,0,pi).integrate(v,0,2*pi)
show('Int = ', Int)
parametric_plot3d([2*sin(u)*cos(v),2*sin(u)*sin(v),2*cos(u)],
                  (u,0,pi),(v,0,2*pi),opacity=0.5)
\end{verbatim}
\end{programcode}
The output is shown in Figure \ref{SphereIntx2}
\begin{figure}[H]
\sidecaption[t]
\includegraphics[scale=.65]{chapters/ch10VectorFields/SphereIntx2}
\caption{Integral of $x^2$ over a sphere}
\label{SphereIntx2}       % Give a unique label
\end{figure}

As an application of such surface integrals, let's consider a cone-shaped surface lamina S given by $z=g(x,y)=4-2\sqrt{x^2+y^2}$ as shown in the Figure below.   At each point on S, the density $\rho = Kr$ is proportional to the radial distance to the z-axis.  Find the mass $m$ of the lamina.   To find intersection with $xy$\-plane, set $z=4-2\sqrt{x^2+y^2}=0$, so $\sqrt{x^2+y^2}=2 \ \Rightarrow \ x^2+y^2=4$ \- circle of radius 2.   
$$g_x(x,y) = -2\frac{2x}{2\sqrt{x^2+y^2}}=-\frac{2x}{\sqrt{x^2+y^2}}$$
$$ g_y(x,y) = -2\frac{2y}{2\sqrt{x^2+y^2}}=-\frac{2y}{\sqrt{x^2+y^2}} \ \ \ \Rightarrow$$
$$ \sqrt{1+g_x^2+g_y^2} = \sqrt{ 1+\frac{4x^2}{x^2+y^2} +\frac{4y^2}{x^2+y^2} }  = \sqrt{5} $$
Then mass is
$$m = \int\int \rho dS = \int_0^{2\pi}\int_0^2 Kr  \sqrt{5} r dr d\theta =  2\pi  \sqrt{5} K \left( \frac{r^3}{3} \right)_0^2 =  \frac{16K\pi \sqrt{5} }{3} $$


\begin{programcode}{SageMath Code and Output:}
\begin{verbatim}
reset(); var('x,y,z,r,K');  
g(x,y) = 4-2*sqrt(x^2+y^2);
gx(x,y) = derivative(g(x,y),x)
gy(x,y) = derivative(g(x,y),y)
dS(x,y) = sqrt(1+gx(x,y)^2+gy(x,y)^2).full_simplify()
show('gx(x,y) = ',gx(x,y),' gy(x,y) = ',
     gy(x,y), ' dS = ', dS(x,y))
f(r) =K*r; 
show('integrand = f(r)*r*sqrt(5) = ', f(r)*r*sqrt(5))
m = ( f(r)*r*sqrt(5) ).integrate(r,0,2).integrate(y,0,2*pi)
show('m = ', m)

show(parametric_plot3d([x,y,g(x,y)],
                  (x,-2,2),(y,-2,2),opacity=0.5))

\end{verbatim}
\end{programcode}
The output is shown in Figure \ref{SurfIntCone}
\begin{figure}[H]
\sidecaption[t]
\includegraphics[scale=.65]{chapters/ch10VectorFields/SurfIntCone}
\caption{Mass of a cone-shaped lamina}
\label{SurfIntCone}       % Give a unique label
\end{figure}


Let's also consider a flux integral of a vector field.   Let S be the portion of the paraboloid $z = g(x, y) = 9 – x^2 – y^2$ lying above the xy-plane, oriented by an upward unit normal vector.   A fluid of constant density $\rho=K$ is flowing through the surface S according to the vector field $F(x, y, z) = <x,y,3z>$.   Find the rate of mass flow through S: $\int\int \rho\mathbf{F}d\mathbf{S}$.   
 $$\int\int \rho\mathbf{F}d\mathbf{S} = \int\int \rho\mathbf{F}\cdot \mathbf{N}dS =  \int\int \rho\mathbf{F}\cdot \frac{\mathbf{r}_u \times \mathbf{r}_v}{||\mathbf{r}_u \times \mathbf{r}_v||} ||\mathbf{r}_u \times \mathbf{r}_v||dA = \int\int \rho\mathbf{F}\cdot (\mathbf{r}_u \times \mathbf{r}_v)dA$$
In this case, the surface is given explicitly: $\mathbf{r} = \left<  x,y,g(x,y) \right>$
$$ \mathbf{n} = \mathbf{r}_u \times \mathbf{r}_v = \begin{pmatrix}
\mathbf{i} & \mathbf{j} & \mathbf{k}\\
1   & 0 &     g_x(x,y) \\
0   & 1 &   g_y(x,y)
\end{pmatrix}  = \left<  -g_x(x,y), -g_y(x,y),1 \right>$$
In our case  $\mathbf{n} = \mathbf{r}_u \times \mathbf{r}_v = \left<  2x, 2y,1 \right>$, so 
$$\mathbf{F}d\mathbf{S} = \left<  x, y, 3z \right> \cdot \left<  2x, 2y,1 \right>=2x^2+2y^2+3z=2x^2+2y^2+3(9-x^2-y^2)=27- x^2-y^2$$
Set  $z = g(x, y) = 9 – x^2 – y^2 = 0$, so $x^2+y^2=9$ - circle of radius 3. Then
 $$\int\int \rho\mathbf{F}d\mathbf{S} = \int_0^3 \int_0^{2\pi} K(27-r^2)r dr d\theta = 2\pi K \left( 27\frac{r^2}{2} - \frac{r^4}{4} \right)_0^{3} =\frac{405\pi K}{2} $$

\begin{programcode}{SageMath Code and Output:}
\begin{verbatim}
reset(); var('x,y,z,r,t,K');  
F(x,y,z) = vector([x,y,3*z])
g(x,y) = 9-(x^2+y^2);
gx(x,y) = derivative(g(x,y),x)
gy(x,y) = derivative(g(x,y),y)
show('gx(x,y) = ',gx(x,y),' gy(x,y) = ', gy(x,y))
dS(x,y) = vector([-gx(x,y),-gy(x,y),1])
FdS = F(x,y,g(x,y)).dot_product(dS(x,y))
show('FdS = ',FdS)
Flux = ( K*(27-r^2)*r ).integrate(r,0,3).integrate(t,0,2*pi)
show('Flux = ',Flux)
parametric_plot3d([x,y,g(x,y)],
                  (x,-2,2),(y,-2,2),opacity=0.5)

\end{verbatim}
\end{programcode}
The output is shown in Figure \ref{FluxInteg}
\begin{figure}[H]
\sidecaption[t]
\includegraphics[scale=.65]{chapters/ch10VectorFields/FluxInteg}
\caption{Flux integral of constant density fluid}
\label{FluxInteg}       % Give a unique label
\end{figure}








\section{Divergence Theorem}

The Divergence Theorem states that for a solid region $Q$ bounded by a closed surface $S$ oriented outwards, $\int\int_S \mathbf{F}d\mathbf{S} = \int\int\int div( \mathbf{F})dV$ for any vector field $ \mathbf{F}$ with continuous partial derivatives.   Let's first verify this theorem for a paraboloid $z=g(x,y)=9-(x^2+y^2)$ above $xy$-plane and vector field $\mathbf{F} = \left< x,y,x^2+y^2+z^2 \right>$.    The closed surface $S$ consists of the upper paraboloid $S_1$ and the bottom flat disk $S_2$ of radius $3$.

First, for $S_1$: $g(x,y) = 9-x^2-y^2; \  \ \Rightarrow \ g_x = -2x; \ \ g_y=-2y \  \ \Rightarrow \ \mathbf{n} = \mathbf{r}_u \times \mathbf{r}_v = \left<  -g_x, -g_y,1 \right>= \left<  2x, 2y,1 \right>$, so
$$\mathbf{F}d\mathbf{S} = \left<  x, y, x^2+y^2+z^2 \right> \cdot \left<  2x, 2y,1 \right>= 2x^2+2y^2+x^2+y^2+z^2=$$
$$ 3x^2+3y^2+(9-x^2-y^2)^2=3r^2+(9-r^2)^2$$
 $$\int\int_{S_1}\mathbf{F}d\mathbf{S} = \int_0^3 \int_0^{2\pi} (3r^2+(9-r^2)^2)r dr d\theta = 2\pi  \int_0^3 (81 -15r^2 + r^4)r dr  = $$
 $$ =  2\pi \left( 81 \frac{r^2}{2} -15 \frac{r^4}{4} + \frac{r^6}{6} \right)_0^3 = \frac{729 \pi}{2}$$

Next, for $S_2$, note that it must have a normal vector directed outwards, which is down in this case. Then for $g(x,y) =0; \  \ \Rightarrow \ g_x = 0; \ \ g_y=0 \  \ \Rightarrow \ \mathbf{n} = \mathbf{r}_u \times \mathbf{r}_v = \left<  g_x,g_y,-1 \right>= \left<  0, 0, -1 \right>$, so
$$\mathbf{F}d\mathbf{S} = \left<  x, y, x^2+y^2+z^2 \right> \cdot \left<  0, 0,-1 \right>= -x^2-y^2=-r^2$$
$$\int\int_{S_1}\mathbf{F}d\mathbf{S} = \int_0^3 \int_0^{2\pi} (-r^2)r dr d\theta = 2\pi  \int_0^3 (-r^3) dr  =  -2\pi \left( \frac{r^4}{4}\right)_0^3 = - \frac{81 \pi}{2}$$
Then overall 
$$\int\int_{S}\mathbf{F}d\mathbf{S} = \int\int_{S_1}\mathbf{F}d\mathbf{S}+\int\int_{S_2}\mathbf{F}d\mathbf{S}=\frac{729 \pi}{2} - \frac{81 \pi}{2} = 324\pi$$
Now, let's consider the divergence integral.   First of all, the divergence of the vector field is $div(\mathbf{F}) = 1+1+2z=2+2z$.  Then
$$ \int\int\int div(\mathbf{F}) dV = \int_0^{2\pi}\int_0^3 \int_0^{9-r^2} 2+2z dz r dr d\theta = 2\pi \int_0^3 \left(2z + z^2 \right)_0^{9-r^2} rdr =$$
$$  2\pi \int_0^3 \left(2(9-r^2) + (9-r^2)^2 \right) rdr =  2\pi \int_0^3 \left(  99r -20r^3 + r^5 \right) rdr =$$
$$ 2\pi \left( 99\frac{r^2}{2} - 20\frac{r^4}{4}+\frac{r^6}{6} \right)_0^3 =324 \pi $$So, we confirmed that they are the same.   The SageMath cells below compute these integrals.  We begin with the flux over the curved parabolic surface:

\begin{programcode}{SageMath Code and Output:}
\begin{verbatim}
reset(); var('x,y,z,r,t,');  
F(x,y,z) = vector([x,y,x^2+y^2+z^2])
g(x,y) = 9-(x^2+y^2);
gx(x,y) = derivative(g(x,y),x)
gy(x,y) = derivative(g(x,y),y)
show('gx(x,y) = ',gx(x,y),' gy(x,y) = ', gy(x,y))
dS(x,y) = vector([-gx(x,y),-gy(x,y),1])
FdS = F(x,y,g(x,y)).dot_product(dS(x,y))
show('FdS = ',FdS)
Flux1 = ( ((9-r^2)^2 + 3*r^2)*r ).
     integrate(r,0,3).integrate(t,0,2*pi)
show('Flux1 = ',Flux1)
plot3d(g(x,y),(x,-3,3),(y,-3,3))
parametric_plot3d([r*cos(t),r*sin(t),9-r^2],(r,0,3),(t,0,2*pi))

\end{verbatim}
\end{programcode}
The output is shown in Figure \ref{DivTheorFlux1}
\begin{figure}[H]
\sidecaption[t]
\includegraphics[scale=.65]{chapters/ch10VectorFields/DivTheorFlux1}
\caption{Flux integral over paraboloid part of the solid for the Divergence Theorem}
\label{DivTheorFlux1}       % Give a unique label
\end{figure}

Let's now find flux over the $xy$-plane, the flat part of the solid.

\begin{programcode}{SageMath Code and Output:}
\begin{verbatim}
g(x,y) = 0;
gx(x,y) = derivative(g(x,y),x)
gy(x,y) = derivative(g(x,y),y)
show('gx(x,y) = ',gx(x,y),' gy(x,y) = ', gy(x,y))
dS(x,y) = -vector([-gx(x,y),-gy(x,y),1])
FdS = F(x,y,g(x,y)).dot_product(dS(x,y))
show('FdS = ',FdS)
Flux2 = ( (-r^2)*r ).integrate(r,0,3).integrate(t,0,2*pi)
show('Flux2 = ',Flux2)

\end{verbatim}
\end{programcode}
The output is shown in Figure \ref{DivTheorFlux2}
\begin{figure}[H]
\sidecaption[t]
\includegraphics[scale=.65]{chapters/ch10VectorFields/DivTheorFlux2}
\caption{Flux integral over flat xy-plane part of the solid for the Divergence Theorem}
\label{DivTheorFlux2}       % Give a unique label
\end{figure}

Finally, let's compare it to the volume integral of the divergence.

\begin{programcode}{SageMath Code and Output:}
\begin{verbatim}
divF(x,y,z) = F.div()(x,y,z)
show('divF = ',divF(x,y,z))
show('Int DivF = ',(divF(x,y,z)*r).integrate(z,0,9-r^2).
                      integrate(r,0,3).integrate(t,0,2*pi))
show('Flux1 + Flux2 = ', Flux1 + Flux2)

\end{verbatim}
\end{programcode}
The output is shown in Figure \ref{DivIntTheor}
\begin{figure}[H]
\sidecaption[t]
\includegraphics[scale=.65]{chapters/ch10VectorFields/DivIntTheor}
\caption{Flux integral over flat xy-plane part of the solid for the Divergence Theorem}
\label{DivIntTheor}       % Give a unique label
\end{figure}

Let's consider a surface integral of a complex vector field 
$$\mathbf{F} = \left< 3x-sin(y^2),2xy+e^{z^2},tan(xy) \right>$$
over a closed surface bounded by cylinder $x^2+y^2=9$, $xy$-plane, and slanted plane $x+z=4$.   The direct surface integral would be impossible to take, however, the divergence theorem $\int\int_S \mathbf{F}d\mathbf{S} = \int\int\int div( \mathbf{F})dV$ turns 3 surface integrals into one volume integral of much simpler scalar field  $div(\mathbf{F}) = 3+2x$.
$$ \int\int\int div(\mathbf{F}) dV = \int_0^{2\pi}\int_0^3 \int_0^{4-r\cos(t)} (3+2r \cos(t)) dz r dr d\theta = $$
$$ \int_0^{2\pi} \int_0^3  (3+2r \cos(t)) ( 4-r \cos(t))  r dr d\theta = $$
$$ \int_0^{2\pi} \int_0^3  (12 +5r \cos(t) - 2r^2 \cos^2(t))  r dr d\theta =$$
$$  \int_0^{2\pi} \left( 6r^2 +5\frac{r^3}{3}\cos(t) - \frac{r^4}{4}(1+\cos(2t)) \right)_0^3 d\theta = \frac{135\pi}{2} $$
Note that $\int_0^{2\pi} \cos(t) dt = \int_0^{2\pi} \cos(2t) dt =0$ because of the periodicity.
The SageMath cell below illustrates this computation as well as the surfaces.

\begin{programcode}{SageMath Code and Output:}
\begin{verbatim}
reset(); var('x,y,z,r,t,');  
F(x,y,z) = vector([3*x-sin(y^2),2*x*y+e^(z^2),tan(x*y)])
divF(x,y,z) = F.div()(x,y,z)
show('divF = ',divF(x,y,z))

show('Int DivF = ',((2*r*cos(t)+3)*r).integrate(z,0,4-r*cos(t)).
                      integrate(r,0,3).integrate(t,0,2*pi))
p1 = implicit_plot3d(x^2+y^2==9,(x,-3,3),(y,-3,3),(z,0,4))
p2 = plot3d(4-x,(x,-3,3),(y,-3,3),opacity=0.5,color='red')
(p1+p2).show()

\end{verbatim}
\end{programcode}
The output is shown in Figure \ref{DivTheorCylinderCut}
\begin{figure}[H]
\sidecaption[t]
\includegraphics[scale=.65]{chapters/ch10VectorFields/DivTheorCylinderCut}
\caption{Application of the Divergence Theorem on a cylinder and slanted plane}
\label{DivTheorCylinderCut}       % Give a unique label
\end{figure}


Another example where the divergence theorem simplifies computations enormously is the surface integral of the vector field  $\mathbf{F} = \left< x^3-sin(z^2),y^3+e^{z^2},z^3-7xy \right>$ over the sphere of radius 3.    Once again, the direct surface integration is out of question for such a complex vector field, especially over a sphere where the parametrization derivatives cross product is rather complex itself (see previous section).  However, the divergence of our vector field is much simpler: $div(\mathbf{F}) = 3x^2+3y^2+3z^2=3\rho^2$, where $\rho$ is the distance from the origin in spherical coordinates.
$$ \int\int_S \mathbf{F} d \mathbf{S} = \int\int\int div(\mathbf{F}) dV = \int_0^{2\pi}\int_0^{\pi} \int_0^3 3\rho^2 \rho^2 sin(\phi) d\rho d\phi d\theta = $$
$$ 2\pi \left( -cos(\phi) \right)_0^{\pi} 3 \left( \frac{\rho^5}{5} \right)_0^3 = \frac{2916\pi}{5}$$
The SageMath cell below illustrates this computation as well as the surface.

\begin{programcode}{SageMath Code and Output:}
\begin{verbatim}
reset(); var('x,y,z,r,t,');  
F(x,y,z) = vector([x^3-sin(z^2),y^3+e^(z^2),z^3-7*x*y])
divF(x,y,z) = F.div()(x,y,z)
show('divF = ',divF(x,y,z))

var('rho,phi,theta')
show('Int DivF = ',((3*rho^2)*rho^2*sin(phi)).integrate(rho,0,3).
                      integrate(phi,0,pi).integrate(theta,0,2*pi))

implicit_plot3d(x^2+y^2+z^2==9,(x,-3,3),(y,-3,3),(z,-3,3))

\end{verbatim}
\end{programcode}
The output is shown in Figure \ref{DivTheorSphere.png}
\begin{figure}[H]
\sidecaption[t]
\includegraphics[scale=.65]{chapters/ch10VectorFields/DivTheorSphere.png}
\caption{Application of the Divergence Theorem on a sphere}
\label{DivTheorSphere.png}       % Give a unique label
\end{figure}


The Divergence Theorem was stated for a simple solid region $Q$ bounded by a closed surface $S$, but it is also valid for regions that are the finite unions of simple solid regions.   For example, let Q be the solid bounded by the closed surfaces $S_1$ and $S_2$ \- inner and outer closed surfaces, respectively.  The solid's surface is the union $S = S_1 \cup S_2$. Note that the normal direction is pointing outside the solid for both $S_1 (\mathbf{-N_1})$ and $S_2  (\mathbf{N_2})$, but for the inner $S_1$ surface $\mathbf{N_1}$ is into the solid.  The simplest plot with two spheres illustrates the idea in the SageMath cell below.

\begin{programcode}{SageMath Code and Output:}
\begin{verbatim}
reset(); var('x,y,z')
p1 = implicit_plot3d(x^2+y^2+z^2==1,(x,-3,3),
             (y,-3,3),(z,-3,3),color='red',opacity = 0.5)
p2 = implicit_plot3d(x^2+y^2+z^2==9,(x,-3,3),
             (y,-3,3),(z,-3,3),opacity = 0.5)
r2 = vector([0,3,0]);  N2 = vector([0,1,0]); 
p3 = plot(N2, start=r2, color='black', thickness = 5) 
r1 = vector([0,1,0]);  N1 = vector([0,-1,0]); 
p4 = plot(N1, start=r1, color='black', thickness = 5) 
r1 = vector([0,1,0]);  N11 = vector([0,1,0]); 
p5 = plot(N11, start=r1, color='black', thickness = 5) 
(p1+p2+p3+p4+p5).show()
\end{verbatim}
\end{programcode}
The output is shown in Figure \ref{DivTheorInnerSurf.png}
\begin{figure}[H]
\sidecaption[t]
\includegraphics[scale=.65]{chapters/ch10VectorFields/DivTheorInnerSurf.png}
\caption{Divergence Theorem illustration with inner surface}
\label{DivTheorInnerSurf.png}       % Give a unique label
\end{figure}

Then
$$\int\int\int div(\mathbf{F}) dV = \int\int_S \mathbf{F}d\mathbf{S} =\int\int_{S_1} \mathbf{F} \cdot (-\mathbf{N_1})d\mathbf{S} + \int\int_{S_2} \mathbf{F} \cdot \mathbf{N_2} d\mathbf{S}=$$
$$-\int\int_{S_1} \mathbf{F} \cdot \mathbf{N_1}d\mathbf{S} +\int\int_{S_2} \mathbf{F} \cdot \mathbf{N_2} d\mathbf{S}$$

Let's now consider an arbitrary surface $S_2$ enclosing electric charge $q$ at the origin.  We can always find a sphere $S_1$ of small enough radius $a$ enclosing the charge at the origin completely inside $S_2$.  The electric field of a point charge is quadratic: $\mathbf{E} = kq\frac{\mathbf{r}}{||\mathbf{r}||^3}=kq \frac{\left< x,y,z \right>}{(x^2+y^2+z^2)^{3/2}}$.
$$div(\mathbf{E} ) = kq\left(  \frac{1}{(x^2+y^2+z^2)^{3/2}} -x \frac{3}{2} \frac{2x}{(x^2+y^2+z^2)^{5/2}}  \right) +$$
$$\left( \frac{1}{(x^2+y^2+z^2)^{3/2}} -y \frac{3}{2} \frac{2y}{(x^2+y^2+z^2)^{5/2}}  \right) + $$
$$\left( \frac{1}{(x^2+y^2+z^2)^{3/2}} - z \frac{3}{2} \frac{2z}{(x^2+y^2+z^2)^{5/2}}   \right) = 0$$
Then
$$\int\int_{S_2} \mathbf{E} \cdot \mathbf{N_2} d\mathbf{S} = \int\int_{S_1} \mathbf{E} \cdot \mathbf{N_1} d\mathbf{S} = 
\int\int_{S_1}   kq\frac{\mathbf{r}}{||\mathbf{r}||^3} \cdot \frac{\mathbf{r}}{||\mathbf{r}||} dS = $$
$$ \int\int_{S_1}   kq\frac{||\mathbf{r}||^2}{||\mathbf{r}||^4} dS = kq\frac{1}{a^2} \int\int_{S_1}    dS = \frac{kq}{a^2}4\pi a^2 = 4\pi k q$$This shows that the flux of electric field of a point charge at the origin is $4\pi k q$ through ANY closed surface that contains it.   Also, note that there is no contradiction that the flux integral of this zero divergence field over a closed sphere containing origin is not zero, because the inverse quadratic field is not continuous at the origin.   I just wanted to mention this result here without any special SageMath calculation as this point is often not made clear in a regular Calculus course.







\section{Stokes Theorem}

The Stokes Theorem states that for an oriented surface $S$ bounded by positively oriented curve $C$:  $\int_C \mathbf{F}d\mathbf{r} = \int\int_S curl( \mathbf{F})d\mathbf{S}$ for any vector field $ \mathbf{F}$ with continuous partial derivatives.   Once again, let's first verify this theorem for a paraboloid $S$: $z=g(x,y)=9-(x^2+y^2)$ above $ xy$-plane and vector field $\mathbf{F} = \left< 5z,3x,2y^2 \right>$.    The surface $S$ boundary is counterclockwise oriented circle $C$ or radius 3 parametrized by $\mathbf{r}(t) = \left< 3\cos(t),3\sin(t),0 \right>$.   The curl of $\mathbf{F}$ is:
$$curl(\mathbf{F}) = \begin{pmatrix}
\mathbf{i}& \mathbf{j} & \mathbf{k}\\
\frac{\partial}{\partial x} & \frac{\partial}{\partial y} & \frac{\partial}{\partial z} \\
5z &   3x &  2y^2
\end{pmatrix} = \left<4y, 5,3 \right> $$

For $S$: $g(x,y) = 9-x^2-y^2; \  \ \Rightarrow \ g_x = -2x; \ \ g_y=-2y \  \ \Rightarrow \ \mathbf{n} = \mathbf{r}_u \times \mathbf{r}_v = \left<  -g_x, -g_y,1 \right>= \left<  2x, 2y,1 \right>$, so
$$curl(\mathbf{F})d\mathbf{S} = \left<  4y, 5, 3 \right> \cdot \left<  2x, 2y,1 \right>dA= (8xy+10y+3)dA $$
 $$\int\int_{S}curl(\mathbf{F})d\mathbf{S}  = \int_0^3 \int_0^{2\pi} (8r^2 \cos(t)\sin(t) +10r \sin(t)+3)r dr d\theta = $$
 $$ 2\pi  3\left( \frac{r^2}{2} \right)_0^3 =27\pi$$
The integrals $\int_0^{2\pi} \sin(t) dt = \int_0^{2\pi} \cos(t)\sin(t) dt =0$ because of the periodicity.

Now, let's consider the line integral.   The derivative of the parametrization $\mathbf{r}(t) =\left< 3\cos(t),3\sin(t),0 \right> $ is  $\mathbf{r}'(t) =\left< -3\sin(t),3\cos(t),0 \right> $.
$$ \int \mathbf{F} d\mathbf{r} =\int \mathbf{F}\cdot \mathbf{r'}(t) dt =  \int_0^{2\pi}  \left< 5z,3x,2y^2 \right> \cdot  \left< -3\sin(t),3\cos(t),0 \right>  dt =$$
$$  \int_0^{2\pi} \left< 5\cdot 0,3 \cdot 3\cos(t),2 \cdot 9\sin^2(t) \right>\cdot  \left< -3\sin(t),3\cos(t),0 \right>  dt=$$
$$   \int_0^{2\pi} 27\cos^2(t) dt=\int_0^{2\pi} 27\frac{1+\cos(2t)}{2} dt= 27\pi$$
So, we confirmed that they are the same.   The SageMath cell below computes these integrals.

\begin{programcode}{SageMath Code and Output:}
\begin{verbatim}
reset(); var('x,y,z,r,t,');  
F(x,y,z) = vector([5*z,3*x,2*y^2])
curlF(x,y,z) = F.curl()(x,y,z)
show('curl of F = ', curlF(x,y,z))
g(x,y) = 9-(x^2+y^2);
gx(x,y) = derivative(g(x,y),x)
gy(x,y) = derivative(g(x,y),y)
show('gx(x,y) = ',gx(x,y),' gy(x,y) = ', gy(x,y))
dS(x,y) = vector([-gx(x,y),-gy(x,y),1])
curlFdS = curlF(x,y,g(x,y)).dot_product(dS(x,y))
show('curlFdS = ',curlFdS)
show('Int curl dS = ',((8*r*cos(t)*r*sin(t)+10*r*sin(t)+3)*r).
                 integrate(r,0,3).integrate(t,0,2*pi))
parametric_plot3d([r*cos(t),r*sin(t),9-r^2],(r,0,3),(t,0,2*pi))
\end{verbatim}
\end{programcode}
The output is shown in Figure \ref{StokesTheor1.png}
\begin{figure}[H]
\sidecaption[t]
\includegraphics[scale=.65]{chapters/ch10VectorFields/StokesTheor1.png}
\caption{Stokes Theorem surface integral of curl}
\label{StokesTheor1.png}       % Give a unique label
\end{figure}

Let's now compute the line integral over the circular border to confirm the Stokes Theorem.

\begin{programcode}{SageMath Code and Output:}
\begin{verbatim}
r(t) = vector([3*cos(t),3*sin(t),0]); 
rp(t) = derivative(r(t),t)
show('r(t) = ',r(t))
show('rp(t) = ',rp(t))
Fdr(t) = F(r(t)[0],r(t)[1],r(t)[2]).dot_product(rp(t))
show('Fdr(t) = ', Fdr(t))
\end{verbatim}
\end{programcode}
The output is shown in Figure \ref{StokesTheor2.png}
\begin{figure}[H]
\sidecaption[t]
\includegraphics[scale=.65]{chapters/ch10VectorFields/StokesTheor2.png}
\caption{Stokes Theorem line integral on the border}
\label{StokesTheor2.png}       % Give a unique label
\end{figure}

Note that the circle of radius 3 in the $xy$-plane also bounds a simple disk in the same plane.   For this surface $g(x,y)=0$, so $\left< -g_x,-g_y,1 \right>=\left< 0,0,1 \right>$, therefore the surface curl integral becomes particularly simple:
 $$\int\int_{S}curl(\mathbf{F})d\mathbf{S}  =\int\int_{S} \left<  4y, 5, 3 \right> \cdot \left<  0, 0,1 \right>dA = 3 \cdot \pi 3^2 = 27\pi $$
This is the same answer that was obtained most easily.

\begin{programcode}{SageMath Code and Output:}
\begin{verbatim}
reset(); var('x,y,z,r,t,');  
F(x,y,z) = vector([5*z,3*x,2*y^2])
curlF(x,y,z) = F.curl()(x,y,z)
show('curl of F = ', curlF(x,y,z))
g(x,y) = 0;
gx(x,y) = derivative(g(x,y),x)
gy(x,y) = derivative(g(x,y),y)
show('gx(x,y) = ',gx(x,y),' gy(x,y) = ', gy(x,y))
dS(x,y) = vector([-gx(x,y),-gy(x,y),1])
curlFdS = curlF(x,y,g(x,y)).dot_product(dS(x,y))
show('curlFdS = ',curlFdS)
show('Int curl dS = ',((3)*r).
                 integrate(r,0,3).integrate(t,0,2*pi))
\end{verbatim}
\end{programcode}
The output is shown in Figure \ref{StokesTheor3.png}
\begin{figure}[H]
\sidecaption[t]
\includegraphics[scale=.65]{chapters/ch10VectorFields/StokesTheor3.png}
\caption{Stokes Theorem surface integral of the curl on a flat disk}
\label{StokesTheor3.png}       % Give a unique label
\end{figure}



Let's consider a line integral of a complex vector field $\mathbf{F} = \left< sin(x^2)-2y^3,2x^3,\tan(z)+e^{z^2} \right>$ over the intersection of the cylinder $x^2+y^2=9$, and a slanted plane $x+z=4$.   The direct line integral would be impossible to take, however, the Stokes theorem $\int_C \mathbf{F}d\mathbf{r} = \int\int_S curl( \mathbf{F})d\mathbf{S}$ has a much simpler curl:
$$curl(\mathbf{F}) = \begin{pmatrix}
\mathbf{i}& \mathbf{j} & \mathbf{k}\\
\frac{\partial}{\partial x} & \frac{\partial}{\partial y} & \frac{\partial}{\partial z} \\
 sin(x^2)-2y^3 &   2x^3  & tan(z)+e^{z^2}
\end{pmatrix} = \left<0, 0,6x^2+6y^2 \right> $$
For $S$: $g(x,y) = 4-x; \  \ \Rightarrow \ g_x = -1; \ \ g_y=0 \  \ \Rightarrow \ \mathbf{n} = \mathbf{r}_u \times \mathbf{r}_v = \left<  -g_x, -g_y,1 \right>= \left<  1, 0,1 \right>$, so
$$curl(\mathbf{F})d\mathbf{S} = \left<  0, 0,6x^2+6y^2 \right> \cdot \left<  1, 0,1 \right>dA= (6x^2+6y^2)dA $$
$$\int_C \mathbf{F}d\mathbf{r} =   \int\int_{S}curl(\mathbf{F})d\mathbf{S}  = \int_0^{2\pi}  \int_0^3 (6r^2)r dr d\theta = 2\pi  6\left( \frac{r^4}{4} \right)_0^3 =243\pi$$

\begin{programcode}{SageMath Code and Output:}
\begin{verbatim}
reset(); var('x,y,z,r,t,');  
F(x,y,z) = vector([sin(x^2)-2*y^3,2*x^3,tan(z)+e^(z^2)])
curlF(x,y,z) = F.curl()(x,y,z)
show('curl of F = ', curlF(x,y,z))
g(x,y) = 4-x;
gx(x,y) = derivative(g(x,y),x)
gy(x,y) = derivative(g(x,y),y)
show('gx(x,y) = ',gx(x,y),' gy(x,y) = ', gy(x,y))
dS(x,y) = vector([-gx(x,y),-gy(x,y),1])
curlFdS = curlF(x,y,g(x,y)).dot_product(dS(x,y))
show('curlFdS = ',curlFdS)
show('Int curl dS = ',((6*r^2)*r).
                 integrate(r,0,3).integrate(t,0,2*pi))
p1 = implicit_plot3d(x^2+y^2==9,(x,-3,3),(y,-3,3),(z,0,4))
p2 = plot3d(4-x,(x,-3,3),(y,-3,3),opacity=0.5,color='red')
(p1+p2).show()
\end{verbatim}
\end{programcode}
The output is shown in Figure \ref{StokesCylinderPlane.png}
\begin{figure}[H]
\sidecaption[t]
\includegraphics[scale=.65]{chapters/ch10VectorFields/StokesCylinderPlane.png}
\caption{Stokes' Theorem of cylinder and plane}
\label{StokesCylinderPlane.png}       % Give a unique label
\end{figure}
